\label{chap:ht-bulk-boundary-line}

% Local macros (not in main.tex preamble).
\providecommand{\kk}{\mathbb{k}}
\providecommand{\SCop}{\mathsf{SC}^{\mathrm{ch,top}}}
\providecommand{\cJ}{\mathcal J}
\providecommand{\cL}{\mathcal L}
\providecommand{\PV}{\operatorname{PV}}
\providecommand{\ot}{\otimes}
\providecommand{\BBar}{\operatorname{Bar}}
\providecommand{\Cobar}{\operatorname{\Omega}}
\providecommand{\Spf}{\operatorname{Spf}}
\providecommand{\mc}{\operatorname{MC}}
\providecommand{\op}{\mathrm{op}}
\providecommand{\der}{\mathrm{der}}
\providecommand{\Kline}{K_{\mathrm{line}}}
\providecommand{\intr}{\mathrm{intr}}
\providecommand{\free}{\mathrm{free}}
\providecommand{\eff}{\mathrm{eff}}
\providecommand{\st}{\mathrm{st}}
\providecommand{\ex}{\mathrm{ex}}
\providecommand{\minmod}{\mathrm{min}}
\providecommand{\expmod}{\mathrm{exp}}
\providecommand{\hkR}{\mathrm{HKR}}
\providecommand{\sn}{\mathrm{SN}}
\providecommand{\hSym}{\widehat{\Sym}}
\providecommand{\cA}{\mathcal{A}}
\providecommand{\Mbar}{\overline{\mathcal{M}}}
\providecommand{\cZ}{\mathcal{Z}}
\providecommand{\cW}{\mathcal{W}}
\providecommand{\gAmod}{\mathfrak{g}^{\mathrm{mod}}_\cA}
\providecommand{\MC}{\operatorname{MC}}
\providecommand{\Conv}{\operatorname{Conv}}
\providecommand{\Tr}{\operatorname{Tr}}


\section{The bulk-boundary-line Koszul triangle}
\label{sec:ht-bulk-boundary-line}

\subsection{The slab and its three algebraic faces}
\index{slab geometry!fiber functor}
\index{Drinfeld double!from slab geometry}
\index{bulk-boundary-line!E1 primacy@$E_1$ primacy}

A 3d holomorphic--topological theory $T$ on $\mathbb{C}_z \times \mathbb{R}_t$
produces three classes of operators: bulk local operators
(T18(iii): supported in the interior of the slab; coincide with
the chiral derived centre T18(i) on the chirally Koszul locus),
boundary local operators (supported on a wall), and line operators
(extended along $\mathbb{R}_t$). The naive
identification ``bulk $=$ boundary $=$ lines'' is false. Bulk
operators form a commutative algebra; boundary operators need not
commute; line operators carry a monoidal structure from fusion.
The deficiency is structural: the three operator classes live at
three different categorical levels. What corrects it?

The answer is forced by the geometry of the slab. Place $T$ on
$\mathbb{C}_z \times I$, where $I = [0,1]$ is an interval with
a boundary condition $D$ on $\{0\}$ and a transverse boundary
condition $N$ on $\{1\}$. Each face supports a chiral algebra:
$\cA$ on the $D$-wall, $\cA^!$ on the $N$-wall
(with $\cA^!$ read from the Ran-space Verdier dual bar coalgebra by
Volume~I, Theorem~A).
The slab $D \,|\, T \,|\, N$ has \emph{two} boundary
components \cite{Dimofte25}. It is a bimodule geometry, not a
Swiss-cheese disk: the Swiss-cheese operad $\SCchtop$ has one
closed interior boundary and one open boundary, while the slab
has two walls of the same type, related by transversality.

A line operator $\ell$ threaded through the slab defines a vector
space
\begin{equation}\label{eq:slab-fiber-functor}
F(\ell) \;:=\; \Hom_T(D,\; \ell,\; N).
\end{equation}
The monoidal structure
$F(\ell \otimes \ell') \simeq F(\ell) \otimes F(\ell')$
follows from factorisation in the topological $\mathbb{R}$-direction:
cutting the slab along a transverse slice decomposes the Hom space.
The functor $F$ is a fiber functor, and the universal algebra
$U = \End(F)$ reconstructing the line category from $F$ is the
\emph{Drinfeld double}
\[
U \;=\; \cA \bowtie \cA^!,
\]
with five structural components, each geometrically forced:
\begin{enumerate}[label=\textup{(\roman*)},leftmargin=1.8em]
\item $\cA$ (the $D$-wall boundary algebra) and $\cA^!$ (the
 $N$-wall Koszul dual) are the two Hopf-algebra factors;
\item the Hopf pairing
 $\langle a, b \rangle$ for $a \in \cA$, $b \in \cA^!$ is the
 hemisphere partition function (the slab with both walls capped);
\item the antipode $S$ comes from orientation reversal of the
 interval (the cup geometry);
\item the coproduct
 $\Delta(a) = a^{(1)} \otimes a^{(2)}$ comes from cutting a line
 into two segments;
\item the $R$-matrix $R(z)$ intertwining $\Delta_z$ and
 $\Delta_z^{\mathrm{op}}$ is the meromorphic braiding from the
 $\mathbb{C}_z$-direction.
\end{enumerate}
The slab forces a bimodule structure because two transverse boundary
conditions produce two independent algebraic inputs, and the
line-operator algebra must intertwine both. The Drinfeld double
$\cA \bowtie \cA^!$ is the unique Hopf algebra receiving compatible
actions from both walls. Complementarity (Theorem~C of Volume~I)
is the algebraic manifestation of the transversality:
$Q_g(\cA) \oplus Q_g(\cA^!) = H^*(\Mbar_g,\, Z(\cA))$; the
deconcatenation coproduct on $B^{\mathrm{ord}}(\cA)$ is the algebraic
avatar of the line-cutting geometry.

\subsection{The corrected triangle}

The slab geometry produces three vertices; the corrected
identification among them is:
\begin{equation}\label{eq:global-corrected-triangle}
\Abulk
\;\simeq\;
\Zder(\Bbound)
\;\simeq\;
\Zder(\cC_{\mathrm{line}})
\;\simeq\;
\HH^\bullet_{\mathrm{ch}}(\cA^!_{\mathrm{line}}),
\qquad
\cC_{\mathrm{line}}\simeq \cA^!_{\mathrm{line}}\text{-mod}
\quad
\text{\small(on the chirally Koszul locus, under compact-generation
and derived-center hypotheses)}.
\end{equation}
Here $\cC_\partial$ denotes the open-sector factorization category
(Definition~\ref{def:oc-factorization-category}); $B_\partial =
\End_{\cC_\partial}(b)$ is the boundary algebra for a chosen vacuum
$b$; and on the chirally Koszul locus
the line category $\cC_{\mathrm{line}}$ is modeled by
$\cA^!_{\mathrm{line}}\text{-}\mathbf{mod}$
(Theorem~\ref{thm:lines_as_modules}). The center descriptions are
Morita invariant
(Corollary~\ref{cor:bulk-derived-center-categorical}). The
unconditional theorem surface is weaker: for any logarithmic
$\SCchtop$-algebra, the bulk identifies with chiral Hochschild
cochains of the boundary, while the line-side module model and the
displayed derived-center identifications require the scoped
hypotheses just stated.

The content of the triangle is a correction, not an identification.
The bulk is the \emph{derived center} of the boundary, not the
boundary itself. The line category is modeled by modules for the
open-colour Koszul dual $\cA^!_{\mathrm{line}}$, not for the bulk algebra.
This is the 3d holomorphic--topological analogue of the strongest
open/closed statements in 2d TFT and the B-model.

In the $E_1$ primacy framework
(Volume~I, Theorem~\ref*{V1-thm:e1-primacy}), the line-operator
vertex $\cC_{\mathrm{line}}$ is the primitive datum: it carries the
full ordered content ($r$-matrix, KZ associator, higher Yangian
coherences). The bulk vertex
$\cZ^{\mathrm{der}}_{\mathrm{ch}}(\cA)$ is its derived center (a
Morita-invariant $\Sigma_n$-coinvariant), and the boundary vertex
$\cA = \mathrm{End}(b)$ is a chart. The triangle descends from the
three-sector decomposition of the $\SCchtop$ convolution algebra
(Volume~I, Corollary~\ref*{V1-cor:convolution-factorization}):
closed factor $\mathfrak{g}^{\mathrm{mod}}$ (bulk/modular), open
factor $\mathfrak{g}^{\mathbb{R}}$ (line/$r$-matrix), mixed sector
(bulk-to-boundary brace evaluation).

Three principles discipline the constructions that follow:
(P1)~intrinsic formulations before complement-dependent
presentations; (P2)~derived or $A_\infty$ models before
cohomology-level data; (P3)~local exact theorems before global
categorical comparisons.

\begin{remark}[Homotopy invariance of the triangle]
\label{rem:bbl-e1-primacy}
\label{rem:triangle-homotopy-first}
\index{bulk-boundary-line!homotopy first}
Each vertex of the corrected triangle is represented by a chosen
strict model (a dg or chiral algebra, an $A_\infty$ boundary
algebra, a dg category of modules), but the triangle itself is
invariant only up to derived/Morita and $L_\infty$ equivalence.
Higher operations from \cite{GKW24} and the line-operator
$A_\infty$ package of \cite{DNP25} fit the same triangle
(Convention~\ref{conv:vol2-strict-models}).
\end{remark}

\begin{remark}[Two-color refinement]
\label{rem:triangle-two-colors}
\index{bulk-boundary-line!two-color refinement}
Closed-color Hochschild theory identifies the bulk:
$A_{\mathrm{bulk}} \simeq C^\bullet_{\mathrm{ch}}(\Bbound,\Bbound)$
(Theorem~\ref{thm:bulk_hochschild}).
On the chirally Koszul locus, open-color Koszul duality identifies
line operators:
$\mathcal{C}_{\mathrm{line}} \simeq
\cA^!_{\mathrm{line}}\text{-mod}$
(Theorem~\ref{thm:lines_as_modules}).
The $R$-matrix lives in the interaction of both colors: it is a
mixed-color datum in the holomorphic weight filtration on $\SCchtop$
(Theorem~\ref{thm:spectral-parameter-from-closed-color}).
On that locus, $\cA^!_{\mathrm{line}}$ carries the dg-shifted
Yangian package (Theorem~\ref{thm:yangian-recognition}), and the full
two-colour Koszul duality involution returns $\cA$
(Theorem~\ref{thm:duality-involution}).

The $\mathsf{E}_1$ modular convolution algebra
$\mathfrak{g}_\cA^{\mathrm{mod},E_1}$
of Vol~I (\S\ref{sec:e1-holographic-bridge}) controls the
open-color structure: the Maurer--Cartan element
$\Theta_\cA^{E_1}$ specialises at degree $2$ to the spectral
kernel $r(z)$, at degree $3$ to the Drinfeld associator, and its
$\Sigma_r$-coinvariant recovers the closed-color shadow obstruction
tower $\kappa, C, Q, \ldots$ of the bulk algebra.
\end{remark}

\begin{remark}[Steinberg presentation]
\label{rem:steinberg-triangle}
\index{bulk-boundary-line!Steinberg presentation}
The three vertices live on three different spaces related to the
Lagrangian inclusion $\cL \hookrightarrow \cM$ and its
self-intersection
$\mathfrak{S}_b = \cL \times_{\cM} \cL$:
\begin{itemize}
\item Bulk:
 $A_{\mathrm{bulk}} \simeq \cO(T^\ast[-1]\cL)$, the algebra of
 polyvector fields on the Lagrangian (equivalently, the Hochschild
 cochains of $\cO(\cL)$). The restriction
 $\cO(\cM)|_{\cL} = \cO(\cL)$ is the boundary, not the bulk; the
 bulk retains the full $(-1)$-shifted cotangent fibre.
\item Boundary: $B_\partial = \cO(\cL)$, functions on the
 Lagrangian itself.
\item Lines:
 $\cC_{\mathrm{line}} = \mathfrak{S}_b\text{-mod}$, modules for
 the convolution algebra on the self-intersection.
\end{itemize}
The Lagrangian condition
$T_{\cL}\cM / T\cL \simeq T^*\cL[-1]$ mediates among all three.
The equivalence $Z_{\der}(B_\partial) \simeq A_{\mathrm{bulk}}$
is the HKR theorem: Hochschild cochains of $\cO(\cL)$ recover
polyvector fields on the shifted cotangent
$T^*[-1]\cL \simeq \cM|_{\cL}$.
The equivalence
$\cC_{\mathrm{line}} \simeq \cA^!_{\mathrm{line}}\text{-mod}$ is the
convolution-module model for the self-intersection algebra.
Both equivalences require the Lagrangian condition, which makes the
normal bundle self-dual and the Steinberg algebra Koszul (with
additional geometric input from purity of the relevant IC~sheaves).
\end{remark}

\begin{remark}[The slab fiber functor and the Drinfeld double]
\label{rem:slab-fiber-functor}
The slab geometry of the preceding subsection realises the
Drinfeld double $\cA \bowtie \cA^!$ as the universal algebra of
line operators \cite{Dimofte25}. The slab is a bimodule with two
boundary components ($D$-wall and $N$-wall), not an $\SCchtop$
disk (which has one open boundary and one closed interior).
This remark records the translation to the monograph's algebraic
framework.
The open/closed MC element
$\Theta^{\mathrm{oc}} = \Theta_A + \sum \mu^{M_j}$ has three
algebraic faces: the universal bulk
$\cZ^{\mathrm{der}}_{\mathrm{ch}}(\cA)$ (chiral derived center),
the boundary algebra $\cA$ (open-color chart), and the line-sector
operations on $\cC_{\mathrm{line}}$ (open-color face). All three
descend from the slab geometry: the $D$-wall gives the boundary,
the Drinfeld double gives the lines, and the derived center of
the boundary gives the bulk.
\end{remark}

\begin{remark}[Dimofte interface generalization]
\label{rem:dimofte-interface-generalization}
\index{slab geometry!interface generalization}
\index{relative Hopf algebra!interface to free theory}
\index{relative bar complex!over Heisenberg}
The slab picture presupposes that the bulk theory $T$ admits
transverse boundary conditions $D$ and $N$. Many 3d
holomorphic--topological theories carry no topological boundary
condition that the line sector can end on \cite{Dimofte25}.
In that situation, one replaces the walls by \emph{interfaces}
from $T$ to a simpler reference theory $\cF$, taken to be a free
theory (Heisenberg $\Heis$, or a supersymmetric refinement) for
which the line-operator category $\cC_\cF$ is understood in closed
form. Let $I_D, I_N \colon T \leftrightsquigarrow \cF$ be two
transverse $T$--$\cF$ interfaces. The slab
$I_D \,|\, T \,|\, I_N$ defines a module over $\cC_\cF$: the
fiber functor refines to
\[
 F_\cF \colon \cC_T \longrightarrow \cC_\cF,
 \qquad
 F_\cF(\ell)
 \;:=\;
 \Hom_T\!\bigl(I_D,\; \ell,\; I_N\bigr),
\]
landing in $\cC_\cF$ rather than in $\Vect$. Setting
$\cF = \text{trivial theory}$ recovers $\cC_\cF = \Vect$ and the
absolute fiber functor of the slab geometry.

The Hopf algebras $\cA$ and $\cA^!$ become \emph{relative} Hopf
algebras internal to $\cC_\cF$. On the chiral algebra side, the
natural object is not the absolute bar coalgebra
$B^{\mathrm{ord}}(\cA)$ but the \emph{relative bar complex}
$B^{\mathrm{ord}}(\cA)/B^{\mathrm{ord}}(\Heis)$, computed by
tensoring the Heisenberg bar over its own deconcatenation coaction.
Relative bar-cobar duality (Volume~I, Theorem~A, applied in the
category $\Heis\text{-}\mathrm{comod}$) produces a relative Koszul
dual $\cA^!_{/\Heis}$, and complementarity (Theorem~C) holds in
$\cC_\cF$ rather than in graded vector spaces. The degeneration
$\cF \to \text{trivial}$, $\Heis \to \mathbf{k}$ collapses
$B^{\mathrm{ord}}(\cA)/B^{\mathrm{ord}}(\mathbf{k})
\simeq B^{\mathrm{ord}}(\cA)$ and recovers the absolute framework.
The interface generalization is the unique extension of the slab
picture to theories whose line sector cannot be cut by a wall but
only by an interface to a free model.
\end{remark}

%%%%%%%%%%%%%%%%%%%%%%%%%%%%%%%%%%%%%%%%%%%%%%%%%%%%%%%%%%%%%%%%%%
%% STAGE 1: THE LOCAL ONE-COLOR THEOREMS
%%%%%%%%%%%%%%%%%%%%%%%%%%%%%%%%%%%%%%%%%%%%%%%%%%%%%%%%%%%%%%%%%%

\section{Stage 1: the local one-color theorems}
\label{sec:stage1-local-theorems}

The four-stage open/closed programme
(Concordance \S\ref{V1-par:concordance-four-stage-architecture})
requires, as its first stage, four purely local theorems that
do not use any global factorization or modularity.
These theorems operate at the level of a \emph{single} chiral
algebra~$\cA$, its bar coalgebra~$\barB(\cA)$,
and the Koszul dual~$\cA^!$ (when $\cA$ is chirally Koszul).
Their proofs use only:
the universal property of the bar construction
(Volume~I, Theorem~A),
bar-cobar inversion on the Koszul locus
(Volume~I, Theorem~B),
the FM compactification and Stokes calculus
(Theorem~\ref{thm:FM-calculus}),
and the homotopy-Koszulity of $\SCchtop$
(Theorem~\ref{thm:homotopy-Koszul}).

\subsection{Theorem S1: MC couplings versus $\Perf(\cA^!)$}
\label{subsec:mc-couplings-perf}

The MC coupling category $\operatorname{Line}(\cA)$
(Definition~\ref{def:mc-coupling-category}) is defined for any
$A_\infty$-chiral algebra without hypotheses. The following
theorem identifies it on the Koszul locus with the algebraic
category of perfect modules over the Koszul dual.

\begin{theorem}[MC couplings $\simeq$ $\Perf(\cA^!)$;
 \ClaimStatusProvedHere{}
 on the chirally Koszul locus]
\label{thm:mc-couplings-perf}
\index{MC coupling category!Perf equivalence|textbf}
Let $\cA$ be a chirally Koszul $A_\infty$-chiral algebra. Then:
\begin{enumerate}[label=\textup{(\roman*)}]
\item The bar construction $\barB(\cA)$ is the universal twisting
 coalgebra: for any conilpotent dg coalgebra~$C$, the set of
 twisting morphisms $\tau \colon C \to \cA$ is in natural bijection
 with coalgebra maps $C \to \barB(\cA)$.

\item There is a quasi-equivalence of dg-categories
 \[
 \Phi \colon \operatorname{Line}(\cA)
 \;\xrightarrow{\;\sim\;}\;
 \Perf(\cA^!),
 \qquad
 (V, \mu)
 \;\longmapsto\;
 (V,\, \rho = (\mathrm{id}_V \otimes \pi)(\mu)),
 \]
 where $\pi \colon \cA \twoheadrightarrow \cA^!$ is the
 bar-cohomology projection.

\item The inverse quasi-equivalence
 $\Psi \colon \Perf(\cA^!) \to \operatorname{Line}(\cA)$
 sends an $\cA^!$-module $(V, \rho)$ to the MC coupling
 $\mu = (\rho \otimes \mathrm{id})(\tau_{\mathrm{univ}})$,
 where $\tau_{\mathrm{univ}} \in \cA^! \otimes \cA$ is the
 universal twisting morphism.
\end{enumerate}
\end{theorem}

\begin{proof}
Part~(i) is the universal property of the bar construction
(Volume~I, Theorem~A, applied to chiral algebras via the
bridge theorem, Theorem~\ref{thm:recognition-SC}). The bar
complex $\barB(\cA)$, equipped with its codifferential,
corepresents the functor $C \mapsto \operatorname{Tw}(C, \cA)$
on the category of conilpotent dg coalgebras.

For part~(ii), on the Koszul locus the bar cohomology
$H^\bullet(\barB(\cA))$ is concentrated in bar degree~$1$ and
equals $\cA^! := (H^\bullet(\barB(\cA)))^\vee$. Bar-cobar
inversion (Volume~I, Theorem~B) gives a quasi-isomorphism
$\Omega(\barB(\cA)) \xrightarrow{\sim} \cA$. The universal
twisting morphism $\tau_{\mathrm{univ}} \in \cA^! \otimes \cA$
mediates between MC elements and module structures:
an MC coupling $\mu \in \End(V) \widehat\otimes \cA$ projects
along $\pi$ to an $\cA^!$-action $\rho$ on~$V$, and this
projection is a quasi-isomorphism of dg-categories by the
bar-cobar Quillen equivalence applied to modules
(Theorem~\ref{thm:homotopy-Koszul}).

Part~(iii) reverses the projection: given $(V, \rho)$,
the formula $\mu = (\rho \otimes \mathrm{id})(\tau_{\mathrm{univ}})$
produces an MC element because $\tau_{\mathrm{univ}}$ satisfies
the MC equation in $\cA^! \otimes \cA$ (this is the twisting
morphism condition, which is equivalent to the coderivation
$d_{\barB}$ squaring to zero).
\end{proof}

\begin{remark}[Scope and conditionality]
\label{rem:mc-couplings-scope}
Part~(i) is unconditional: it holds for any $A_\infty$-chiral
algebra. Parts~(ii)--(iii) require chiral Koszulness. Off the
Koszul locus, $\operatorname{Line}(\cA)$ is still well-defined
(Remark~\ref{rem:off-koszul-locus}), but the algebraic
identification with $\Perf(\cA^!)$ fails: the bar cohomology may
not be concentrated, and the projection $\pi$ may not exist.
The theorem refines
Theorem~\ref{thm:koszul-comparison-mc} by making the role of
the universal twisting morphism explicit.
\end{remark}

\subsection{Theorem S2: meromorphic tensor dg-category}
\label{subsec:meromorphic-tensor-cat}

The line-operator dg-category carries a meromorphic tensor
structure from fusion at holomorphic separation~$z$.

\begin{theorem}[Meromorphic tensor dg-category;
 \ClaimStatusProvedHere]
\label{thm:meromorphic-tensor-category}
\index{meromorphic tensor category!Stage~1|textbf}
Let $\cA$ be a chirally Koszul $A_\infty$-chiral algebra,
$\cA^!$ its Koszul dual, and $r(z) \in (\cA^!)^{\otimes 2}((z^{-1}))$
the spectral $r$-matrix
\textup{(}Definition~\textup{\ref{def:spectral_R}}\textup{)}.
The triple
$(\operatorname{Line}(\cA),\, \otimes_z,\, \mathbf{1})$
is a meromorphic tensor dg-category:
\begin{enumerate}[label=\textup{(\roman*)}]
\item \emph{Binary fusion.}
 For $M = (V, \mu_V),\; N = (W, \mu_W) \in \operatorname{Line}(\cA)$,
 the meromorphic tensor product $M \otimes_z N$ has underlying
 complex $V \otimes W$ and MC coupling
 $\mu_{M \otimes_z N} = \mu_V \otimes 1 + 1 \otimes \mu_W
 + r(z)$.
 This is a well-defined MC element because $r(z)$ satisfies the
 MC equation in $(\cA^!)^{\otimes 2}$
 \textup{(}Theorem~\textup{\ref{thm:spectral_R_YBE}(i)}\textup{)}.

\item \emph{Unit.}
 The trivial line $\mathbf{1} = (k, 0)$
 \textup{(}from the counit $\varepsilon \colon \cA^! \to k$\textup{)}
 satisfies $\mathbf{1} \otimes_z M \simeq M$ and
 $M \otimes_z \mathbf{1} \simeq T_z M$, where $T_z$ is the
 translation automorphism.

\item \emph{Associativity.}
 $(M \otimes_{z_1} N) \otimes_{z_2} P
 \simeq M \otimes_{z_1+z_2} (N \otimes_{z_2} P)$
 up to coherent quasi-isomorphism, via the coassociativity
 identity
 \[
 \Delta(z_{12}) \circ (\Delta(z_{23}) \otimes \mathrm{id})
 = (\mathrm{id} \otimes \Delta(z_{23})) \circ \Delta(z_{12}+z_{23}),
 \]
 where $\Delta_z \colon \cA^! \to \cA^! \otimes_z \cA^!$
 is the $z$-deformed coproduct from the dg-shifted Yangian
 structure.
\end{enumerate}
\end{theorem}

\begin{proof}
(i) The MC equation for $\mu_{M \otimes_z N}$ reduces to
$d(\mu_V \otimes 1 + 1 \otimes \mu_W + r(z))
+ \tfrac{1}{2}[\mu_V \otimes 1 + 1 \otimes \mu_W + r(z),\,
\mu_V \otimes 1 + 1 \otimes \mu_W + r(z)] = 0$.
The cross terms between $\mu_V$, $\mu_W$ vanish because
the two tensor factors act on different spaces. The remaining
terms are: the MC equations for $\mu_V$ and $\mu_W$ separately
(which hold by hypothesis), and the MC equation for $r(z)$
in $(\cA^!)^{\otimes 2}$ (Theorem~\ref{thm:spectral_R_YBE}(i)).

(ii) When $\mu_V = 0$, the MC coupling reduces to
$\mu_W + r(z)|_{\varepsilon \otimes \mathrm{id}}$. The counit
$\varepsilon$ kills $r(z)$ by the counit axiom of the
dg-shifted Yangian, so $\mu_{\mathbf{1} \otimes_z M} = \mu_M$.
The right-unit statement uses the same argument with the
translation automorphism $T_z$ absorbing the residual spectral
dependence.

(iii) The two parenthesizations produce the MC couplings
\begin{align*}
 \mu_{(M \otimes_{z_1} N) \otimes_{z_2} P}
 &= \mu_V + \mu_W + \mu_P
 + r_{12}(z_1) + r_{(12)3}(z_2), \\
 \mu_{M \otimes_{z_1+z_2} (N \otimes_{z_2} P)}
 &= \mu_V + \mu_W + \mu_P
 + r_{23}(z_2) + r_{1(23)}(z_1{+}z_2).
\end{align*}
These agree by the $A_\infty$ Yang--Baxter relation
(Theorem~\ref{thm:spectral_R_YBE}(ii)),
which is itself a consequence of Stokes' theorem on $\FM_3(\C)$
(Theorem~\ref{thm:YBE}). The coherent higher data come from
higher boundary strata of $\FM_n(\C)$.
\end{proof}

\begin{remark}[Relation to existing statements]
Theorem~\ref{thm:meromorphic-tensor-category} is equivalent to
Theorem~\ref{thm:meromorphic-tensor-dgcat} in the line-operator
chapter. The present formulation isolates it as a Stage~1 theorem,
emphasizing that its proof requires only local data (the $r$-matrix
and its Yang--Baxter equation) and no global factorization.
\end{remark}

\subsection{Theorem S3: meromorphic tensor $\simeq$
 factorization dg-cosheaf}
\label{subsec:meromorphic-factorization}

The meromorphic tensor product~$\otimes_z$ describes
fusion at finite separation. The FM compactification provides the
correct regime for taking $z \to 0$ limits. Passing between the
two presentations is the content of the following equivalence.

\begin{theorem}[Meromorphic tensor $\simeq$ factorization
 dg-cosheaf; \ClaimStatusProvedHere{} on the chirally Koszul locus]
\label{thm:meromorphic-factorization-equiv}
\index{factorization!meromorphic tensor equivalence|textbf}
On the chirally Koszul locus, the meromorphic tensor presentation
$(\operatorname{Line}(\cA), \otimes_z)$
and the factorization dg-cosheaf presentation
$\cC_{\mathrm{op},\partial}$ on ordered boundary configuration
spaces determine each other:
\begin{enumerate}[label=\textup{(\roman*)}]
\item \emph{Tensor to factorization.} On each chamber of
 $\operatorname{Conf}_n^{\mathrm{ord}}(I_p)$, assign the
 external tensor product
 $\cC_J^{\mathrm{op}} \boxtimes \cdots \boxtimes
 \cC_J^{\mathrm{op}}$. On codimension-one collision faces,
 glue by the iterated meromorphic tensor product~$\otimes_z$.
 The codimension-two cocycle condition holds because:
 \textup{(a)}~disjoint collisions commute tautologically;
 \textup{(b)}~triple collisions reduce to the associativity
 of~$\otimes_z$
 \textup{(}Theorem~\textup{\ref{thm:meromorphic-tensor-category}(iii)}\textup{)},
 which is in turn the $A_\infty$ Yang--Baxter relation.

\item \emph{Factorization to tensor.} Given a factorization
 dg-cosheaf $\cC$ on ordered boundary configuration spaces,
 restrict to two-point collision faces to recover
 the binary meromorphic tensor product~$\otimes_z$. Higher
 fusion functors are recovered from multi-point collision
 strata by the factorization structure on each stratum.
\end{enumerate}
Both directions are quasi-inverse: the round-trip
tensor $\to$ factorization $\to$ tensor recovers~$\otimes_z$
by restriction, and factorization $\to$ tensor $\to$
factorization recovers the cosheaf by the universal property
of the FM-type gluing.
\end{theorem}

\begin{proof}
Direction~(i): a constructible dg-cosheaf on the
FM-type compactification of ordered configurations requires a
cocycle condition at every codimension-two face. These faces
are of two types. First, \emph{disjoint collisions}: two
clusters of points collide independently. The gluing data for
disjoint clusters involve independent copies of~$\otimes_z$
acting on different tensor factors, so they commute
tautologically. Second, \emph{nested collisions}: three points
$z_1, z_2, z_3$ collide with two possible parenthesizations,
$((z_1, z_2), z_3)$ vs.\ $(z_1, (z_2, z_3))$. The two
gluings produce $(M_1 \otimes_{z_{12}} M_2) \otimes_{z_{23}} M_3$
and $M_1 \otimes_{z_{12}+z_{23}} (M_2 \otimes_{z_{23}} M_3)$
respectively. These agree by the associativity of~$\otimes_z$
(Theorem~\ref{thm:meromorphic-tensor-category}(iii)). Thus the
cocycle condition is satisfied, and the chamber-wise data extend
to a constructible cosheaf.

Direction~(ii): the factorization axiom
$\cC(U_1 \sqcup U_2) \simeq \cC(U_1) \boxtimes \cC(U_2)$
on disjoint open sets determines the cosheaf on chambers. The
restriction to a codimension-one collision face
$D_{\{i,j\}} \subset \partial \FM_n$ recovers a binary functor
$\cC_J^{\mathrm{op}} \times \cC_J^{\mathrm{op}} \to
\cC_J^{\mathrm{op}}$, which is the meromorphic tensor
product~$\otimes_{z_{ij}}$ by construction of the boundary
residue.

The round-trip identities follow from the two-out-of-three
property: the factorization axiom and the cocycle condition
together determine the cosheaf uniquely from its restriction
to chambers and codimension-one faces.
\end{proof}

\begin{remark}[The FM compactification as collision regulator]
\label{rem:fm-collision-regulator}
The FM compactification~$\FM_n(\C)$ is a blowup along diagonals,
not the complement of the diagonal. Its boundary divisors
$D_S$ ($|S| \ge 2$) parametrize configurations where the
points in~$S$ collide while retaining the information of their
relative positions as an element of $\FM_{|S|}(\C)$.
This blown-up boundary, not the naive $z \to 0$ limit,
provides the correct collision regime:
the meromorphic tensor product~$\otimes_z$ has poles at $z = 0$,
but the FM boundary replaces the pole with a resolved stratum
carrying the OPE data.
\end{remark}

\begin{remark}[Relation to existing statements]
This theorem refines
Theorem~\ref{thm:meromorphic-factorization-equivalence}
in the foundations (\S\ref{subsec:local-to-global-factorization})
by isolating the local content and making the codimension-two
cocycle verification explicit. The global Weiss descent
(Conjecture~\ref{conj:weiss-descent}) is the categorical
boundary analogue, and remains conjectural.
\end{remark}

\subsection{Theorem S4: the boundary algebra as
 $A_\infty$-chiral algebra}
\label{subsec:boundary-ainfty-chiral}

The following theorem constructs the boundary algebra and its
$A_\infty$-chiral structure from a compact generator of the
open-sector factorization category.

\begin{theorem}[Boundary algebra is $A_\infty$-chiral;
 \ClaimStatusProvedHere]
\label{thm:boundary-ainfty-chiral}
\index{boundary algebra!A-infinity chiral|textbf}
Let $\cC$ be an open/closed factorization dg-category on a
tangential log curve $(X, D, \tau)$
\textup{(}Definition~\textup{\ref{def:oc-factorization-category}}\textup{)},
and let $b \in \cC(J)$ be a compact generator for a boundary
interval $J \subset I_p$. Then:
\begin{enumerate}[label=\textup{(\roman*)}]
\item \emph{$A_\infty$-chiral structure.}
 The endomorphism algebra
 \[
 A_b \;:=\; \RHom_{\cC(J)}(b, b)
 \]
 carries an $A_\infty$-chiral algebra structure
 \textup{(}Definition~\textup{\ref{def:ainfty_chiral}}\textup{)}.
 The $k$-th higher product
 \[
 m_k \colon
 C_\ast\bigl(\FM_k(\C)\bigr) \otimes A_b^{\otimes k}
 \;\longrightarrow\;
 A_b((\lambda_1))\cdots((\lambda_{k-1}))
 \]
 arises from restricting the factorization functors of~$\cC$
 to the mixed collision stratum with $k$~closed insertions on
 the endomorphism complex of~$b$.
 The holomorphic variables $\lambda_i = z_i - z_{i+1}$ are the
 OPE spectral parameters from $\FM_k(\C)$; the topological
 ordering from $\Conf_k^{\mathrm{ord}}(\R)$ gives the planar
 $A_\infty$ convention by local constancy
 \textup{(}axiom~\textup{(iv)} of
 Definition~\textup{\ref{def:oc-factorization-category}}\textup{)}.

\item \emph{$A_\infty$ identities.}
 The Stasheff relation
 $\sum_{i+j=n+1} m_i \circ m_j = 0$
 follows from Stokes' theorem on
 $\FM_k(\C) \times \Conf_k^{\mathrm{ord}}(\R)$.
 The codimension-one boundary strata of this product are
 products of lower-dimensional operation spaces, and each
 such stratum corresponds to a composition $m_i \circ m_j$
 with $i + j = k + 1$. The Stokes identity on the total
 boundary
 $\sum_\sigma \int_{\partial_\sigma} = 0$
 is the $A_\infty$ relation.

\item \emph{Morita invariance.}
 A different compact generator $b' \in \cC(J')$ yields an
 $A_\infty$-chiral algebra $A_{b'}$ that is $A_\infty$-Morita
 equivalent to~$A_b$ via the bimodule
 $M_{b,b'} = \RHom_{\cC(J)}(b, b')$.
 The true invariant is the factorization dg-category~$\cC$,
 not a chosen algebra presentation.
\end{enumerate}
\end{theorem}

\begin{proof}
Part~(i):
By local constancy (axiom~(iv) of
Definition~\ref{def:oc-factorization-category}), the dg-category
$\cC(J)$ is independent of the interval $J \subset I_p$ up to
canonical equivalence. The factorization structure of~$\cC$
provides, for each $k$-element configuration
$(z_1, \ldots, z_k) \in \FM_k(\C)$ and ordered open
configuration $(t_1 < \cdots < t_k) \in \Conf_k^{\mathrm{ord}}(\R)$,
a multilinear functor
$\mu_k \colon \cC(J_1) \times \cdots \times \cC(J_k) \to \cC(J)$,
where the $k$~ordered boundary insertions divide~$J$ into $k+1$
subintervals. Restricting $\mu_k$ to
$\Hom_{\cC(J)}(b,b)$ in each slot produces the operation~$m_k$.

The holomorphic parameters $\lambda_i = z_i - z_{i+1}$ arise
from the $\FM_k(\C)$ factor: the operation $m_k$ is a formal
Laurent series in $\lambda_1, \ldots, \lambda_{k-1}$ with
poles from collisions $z_i \to z_{i+1}$.
The topological ordering from $\Conf_k^{\mathrm{ord}}(\R)$
ensures the planar $A_\infty$ convention: inputs are ordered
$a_1, \ldots, a_k$ according to $t_1 < \cdots < t_k$.
Local constancy ensures the result depends only on the
combinatorial ordering, not on positions.
Sesquilinearity follows from translation covariance in the
holomorphic coordinate.

Part~(ii):
The $A_\infty$ identity is Stokes' theorem applied to the
logarithmic weight form $\omega_k^\partial$ on the operation
space $\FM_k(\C) \times \Conf_k^{\mathrm{ord}}(\R)$.
On the interior, $d\omega_k^\partial = 0$ (the weight form is
closed). The codimension-one boundary strata of the product
decompose as
$\bigl(\FM_i(\C) \times \Conf_i^{\mathrm{ord}}(\R)\bigr)
\times
\bigl(\FM_j(\C) \times \Conf_j^{\mathrm{ord}}(\R)\bigr)$
with $i + j = k + 1$. Each such stratum corresponds to a
Stasheff face. The Stokes identity
$\int_{\partial(\FM_k \times \Conf_k)} \omega_k^\partial = 0$
gives the $A_\infty$ relation. The local model axiom~(v) of
Definition~\ref{def:oc-factorization-category} ensures these
boundary identifications are compatible with the
$W(\SCchtop)$-algebra coherence data.

Part~(iii):
Let $b' \in \cC(J)$ be another compact generator. The bimodule
$M_{b,b'} = \RHom_{\cC(J)}(b, b')$ is an $A_\infty$-bimodule
over $(A_b, A_{b'})$. Since both are compact generators,
$M_{b,b'} \otimes^{\mathbf{L}}_{A_{b'}} M_{b,b'}^\vee
\simeq A_b$ and
$M_{b,b'}^\vee \otimes^{\mathbf{L}}_{A_b} M_{b,b'}
\simeq A_{b'}$, which is the Morita equivalence criterion for
compact generators of a dg-category \cite{Kel06,HA}.
\end{proof}

\begin{remark}[Unconditional scope]
\label{rem:boundary-ainfty-unconditional}
Theorem~\ref{thm:boundary-ainfty-chiral} is unconditional:
it requires only the axioms of
Definition~\ref{def:oc-factorization-category} and uses no
Koszulness, formality, or finiteness beyond the compact-generation
hypothesis on~$b$. The chiral Koszulness of~$A_b$ is a further
property that, when it holds, activates the equivalence
$\operatorname{Line}(\cA) \simeq \Perf(\cA^!)$ of
Theorem~\ref{thm:mc-couplings-perf} and the meromorphic tensor
structure of Theorem~\ref{thm:meromorphic-tensor-category}.
\end{remark}

\begin{remark}[Relation to existing content]
\label{rem:stage1-cross-references}
The four Stage~1 theorems consolidate and refine results
already present in scattered locations:
Theorem~\ref{thm:mc-couplings-perf} refines
Theorem~\ref{thm:koszul-comparison-mc};
Theorem~\ref{thm:meromorphic-tensor-category} is equivalent to
Theorem~\ref{thm:meromorphic-tensor-dgcat};
Theorem~\ref{thm:meromorphic-factorization-equiv} refines
Theorem~\ref{thm:meromorphic-factorization-equivalence};
Theorem~\ref{thm:boundary-ainfty-chiral} refines
Theorem~\ref{thm:boundary-algebra-from-vacuum}.
The consolidation serves the four-stage programme: having
all four in one place, with uniform notation and explicit
scope qualifications, is the prerequisite for Stage~2
(open primitive: $\cC_{\mathrm{op}}$ as primitive, Morita
invariance) and Stage~3 (globalization on tangential log curves).
\end{remark}

%%%%%%%%%%%%%%%%%%%%%%%%%%%%%%%%%%%%%%%%%%%%%%%%%%%%%%%%%%%%%%%%%%
%% END OF STAGE 1
%%%%%%%%%%%%%%%%%%%%%%%%%%%%%%%%%%%%%%%%%%%%%%%%%%%%%%%%%%%%%%%%%%

\section{Literature backbone, formal consequences, and status ledger}\label{sec:backbone}

Before constructing anything new, we fix the status of the foundational inputs.

\subsection{What is already in the literature}

The boundary-chiral-algebra paper \cite{CDG20} establishes the cohomology-level bulk algebra, its shifted Poisson bracket, the action of bulk on boundary operators, and the expectation of higher derived operations and of a derived bulk--boundary comparison map. Zeng's 2023 paper \cite{Zeng23} gives strong evidence for the bulk/boundary comparison by identifying boundary chiral algebras with universal defect algebras in families of examples. The higher-operations paper \cite{GKW24} supplies the perturbative technology that constructs higher operations and proves the quadratic identities they satisfy in a broad holomorphic--topological setting. The 2025 line paper \cite{DNP25} models the line category by modules over a Koszul-dual $A_\infty$ algebra and constructs the dg-shifted-Yangian structures organizing their OPE.

These are the literature-backed pillars on which the present development stands.

\subsection{Formal consequences}

Once one has:
\begin{enumerate}[label=(\roman*),leftmargin=1.8em]
\item a rich boundary condition $b$ whose boundary algebra $B_b$ generates the line category;
\item an equivalence $\cC_{\mathrm{line}}\simeq B_b\text{-mod}$ or, more invariantly, a compact generator $M_b\in \cC_{\mathrm{line}}$ with $B_b\simeq \End_{\cC_{\mathrm{line}}}(M_b)$;
\item and a quasi-isomorphism $\beta_{\der}:\Abulk\to \Zder(B_b)$,
\end{enumerate}
the global triangle \eqref{eq:global-corrected-triangle} becomes formal: centers are Morita invariant, so $\Zder(B_b)\simeq \Zder(\cC_{\mathrm{line}})$; if $\cC_{\mathrm{line}}\simeq \cA^!_{\mathrm{line}}\text{-mod}$, then $\Zder(\cC_{\mathrm{line}})\simeq \HH^\bullet_{\mathrm{ch}}(\cA^!_{\mathrm{line}})$.

This part is not new mathematics. It is the clean formal envelope inside which the real work must be done.

\begin{proposition}[Formal reduction of the global triangle; \ClaimStatusProvedHere]\label{prop:formal-global-triangle}
Assume there is a boundary condition $b$ such that:
\begin{enumerate}[label=(\roman*),leftmargin=1.8em]
\item $b$ defines a compact generator $M_b$ of the line category $\cC_{\mathrm{line}}$;
\item $B_b:=\End_{\cC_{\mathrm{line}}}(M_b)$ is a model for the boundary open sector;
\item the derived bulk--boundary map $\beta_{\der}:\Abulk\to \Zder(B_b)$ is a quasi-isomorphism;
\item $\cC_{\mathrm{line}}\simeq \cA^!_{\mathrm{line}}\text{-mod}$ for an open-colour Koszul-dual $A_\infty$ algebra $\cA^!_{\mathrm{line}}$.
\end{enumerate}
Then there are canonical quasi-isomorphisms
\[
\Abulk
\;\simeq\;
\Zder(B_b)
\;\simeq\;
\Zder(\cC_{\mathrm{line}})
\;\simeq\;
\HH^\bullet_{\mathrm{ch}}(\cA^!_{\mathrm{line}}).
\]
\end{proposition}

\begin{proof}
The first equivalence is hypothesis (iii). Since $M_b$ is a compact generator, $\cC_{\mathrm{line}}$ is derived Morita equivalent to $B_b\text{-mod}$. Derived centers are Morita invariant, giving $\Zder(B_b)\simeq \Zder(\cC_{\mathrm{line}})$. Hypothesis (iv) identifies $\cC_{\mathrm{line}}$ with $\cA^!_{\mathrm{line}}\text{-mod}$, and the derived center of the module category is chiral Hochschild cochains of $\cA^!_{\mathrm{line}}$.
\end{proof}

\subsection{What is new in this manuscript}

The novelty of this note lies elsewhere: it strengthens the microlocal algebraic core. We prove, inside a large but exact local sector, that the boundary, line, and bulk vertices can each be constructed explicitly and compared by theorems rather than analogies.

\begin{center}
\renewcommand{\arraystretch}{1.2}
\begin{longtable}{|p{0.24\textwidth}|p{0.23\textwidth}|p{0.43\textwidth}|}
\hline
\textbf{Layer} & \textbf{Status} & \textbf{Content} \\
\hline
Bulk shifted Poisson/chiral algebra; boundary module structure & literature & Costello--Dimofte--Gaiotto \cite{CDG20}. \\
\hline
Higher perturbative operations and quadratic identities & literature & Gaiotto--Kulp--Wu \cite{GKW24}. \\
\hline
Universal defect evidence for the bulk/boundary comparison & literature & Zeng \cite{Zeng23}. \\
\hline
Line category as modules over $\cA^!_{\mathrm{line}}$; dg-shifted Yangian package & literature & Dimofte--Niu--Py \cite{DNP25}. \\
\hline
Formal implication ``rich boundary + bulk/boundary comparison $\Rightarrow$ corrected triangle'' & formal & Proposition~\ref{prop:formal-global-triangle}. \\
\hline
Intrinsic boundary $A_\infty$ algebra from normal Taylor classes & new, proved here & Theorem~\ref{thm:intrinsic-boundary}. \\
\hline
Boundary Thom--Sebastiani and stabilization calculus & new, proved here & Theorems~\ref{thm:TS}, \ref{thm:contractible-stab}, \ref{thm:clifford-stab}, \ref{thm:relative-morse}. \\
\hline
Boundary-linear LG theorem: bulk $=$ shifted cotangent of the derived boundary zero locus & new, proved here & Theorems~\ref{thm:dcrit-shiftedcot}, \ref{thm:boundary-linear-bulk-boundary}. \\
\hline
One-step Jacobi coalgebra and exact pointed line algebra & new, proved here & Theorems~\ref{thm:exact-line}, \ref{thm:mc-vacua}, \ref{thm:polarized-reconstruction}. \\
\hline
Boundary Kuranishi elimination and minimal line algebra & new, proved here & Theorems~\ref{thm:boundary-kuranishi}, \ref{thm:minimal-pointed-line}. \\
\hline
Boundary-linear line Mather--Yau theorem and microlocal classification & new, proved here & Theorem~\ref{thm:line-mather-yau}. \\
\hline
All-theory global Koszul triangle & conjectural & See Section~\ref{sec:global}. \\
\hline
\end{longtable}
\end{center}


\section{First principles: locality, centers, and the correct shape of open/closed/line structures}\label{sec:first-principles}

The boundary and line sectors do not stand beside the bulk as parallel commutative algebras. They sit one categorical level lower. The local operator algebra is bulk. The algebra of endomorphisms of a boundary vacuum or of a line defect is an open-sector presentation. The bridge from open data back to bulk data is the derived center.

This principle is older than the present subject. In 2d open/closed TFT and in the B-model, the closed sector is recovered from the derived center of the open sector. In representation theory, centers are traced not by forgetting categories to algebras but by passing from convolution algebras to inertia geometry. The same architecture reappears here.

\subsection{Why centers must appear}

Suppose $\cC$ is a dg or $A_\infty$ category representing a line sector. Its center is not a subalgebra of a chosen endomorphism algebra; it is the derived endomorphism algebra of the identity functor. If $\cC\simeq A\text{-mod}$, then the center of $\cC$ is Hochschild cochains of $A$:
\[
\Zder(\cC)\simeq \RHom_{\operatorname{Fun}(\cC,\cC)}(\id_{\cC},\id_{\cC})\simeq \HH^\bullet(A).
\]
This is the precise sense in which bulk can be the center of the line sector while lines themselves are controlled by a noncommutative algebra.

\subsection{Why Morita invariance matters}

If a rich boundary condition $b$ generates the line category and $M_b$ denotes the corresponding boundary vacuum object, then $B_b=\End(M_b)$ is only a presentation. Different generators yield different algebras but equivalent module categories. Bulk reconstruction therefore cannot depend on the chosen presentation. Derived center and Hochschild cochains are the correct invariants precisely because they are Morita invariant.

\subsection{Why the exact local sector is decisive}

A global theorem comparing bulk, boundary, and line categories in every holomorphic--topological theory is too large to attack head-on. One needs a local exact sector where the objects admit literal formulas. The boundary-linear Landau--Ginzburg sector plays that role. In that sector:
\begin{itemize}[leftmargin=1.8em]
\item the boundary algebra is an explicit commutative dg algebra;
\item the local pointed line algebra is an explicit $A_\infty$ algebra;
\item the derived center can be computed directly;
\item Kuranishi elimination is exact rather than merely homotopy-theoretic.
\end{itemize}
What follows is the construction of that exact local backbone.


\section{Intrinsic boundary $A_\infty$ algebra from first principles}\label{sec:intrinsic}

Let $V$ be a finite-dimensional vector space, let $L\subset V$ be a linear subspace, and let
\[
Q:=V/L.
\]
Let $I_L\subset \hSym(V^\vee)$ denote the ideal of functions vanishing on $L$, so the completed formal neighborhood of $L$ in $V$ is
\[
\widehat{V}_L:=\Spf\bigl(\hSym(V^\vee)^\wedge_{I_L}\bigr).
\]
Assume $W\in \hSym(V^\vee)^\wedge_{I_L}$ vanishes on $L$, equivalently $W\in I_L$.

The aim is to define a boundary $A_\infty$ algebra depending only on the formal germ of $W$ along $L$, with no choice of complement.

\subsection{Intrinsic normal Taylor classes}

The normal Taylor coefficients of $W$ along $L$ are the classes
\[
\tau_n(W)\in I_L^n/I_L^{n+1}\cong \widehat{\cO}(L)\ot \Sym^n(Q^\vee),
\qquad n\ge 1.
\]
Thus $W$ determines and is determined by the sequence $\{\tau_n(W)\}_{n\ge 1}$.

Choose, only locally for writing formulas, a splitting $V=L\oplus N$. Then $Q\cong N$ and we may write
\[
W(x,y)=\sum_{n\ge 1} W_n(x;y),
\qquad
W_n(x;\lambda y)=\lambda^n W_n(x;y).
\]
Then $\tau_n(W)$ is exactly the class of $W_n(x;y)$.

\subsection{The underlying graded algebra}

Set
\[
B^{\intr,\gr}_{L,W}:=\widehat{\cO}(L)\ot \Lambda(Q),
\]
where the copy of $Q$ sits in cohomological degree $1$. If $q_1,\dots,q_r$ is a basis of $Q$ and $\psi_1,\dots,\psi_r$ the corresponding degree-$1$ generators, then $B^{\intr,\gr}_{L,W}$ is the free completed graded-commutative algebra on $\widehat{\cO}(L)$ and the odd variables $\psi_a$.

The intrinsic Taylor classes of $W$ will produce higher multiderivations on this free algebra.

\subsection{From Taylor classes to multiderivations}

Fix a splitting temporarily and choose coordinates $x_i$ on $L$ and basis vectors $e_a$ of $Q$. Write
\[
\tau_n(W)=\sum_{a_1,\dots,a_n} f_{a_1\cdots a_n}(x)\, e_{a_1}^\vee\cdots e_{a_n}^\vee.
\]
Define an $n$-ary Hochschild cochain $\mu_n$ on $B^{\intr,\gr}_{L,W}$ by
\begin{equation}\label{eq:intrinsic-mu}
\mu_n(a_1,\dots,a_n)
:=
\sum_{a_1,\dots,a_n}
f_{a_1\cdots a_n}(x)\,
(\partial_{\psi_{a_1}}a_1)\cdots(\partial_{\psi_{a_n}}a_n).
\end{equation}
Each $\mu_n$ is a multiderivation of the free graded-commutative product.

The collection $\{\mu_n\}_{n\ge 1}$ is not arbitrary: it comes from the single function $W$, hence satisfies the Stasheff system.

\begin{theorem}[Intrinsic boundary $A_\infty$ algebra; \ClaimStatusProvedHere]\label{thm:intrinsic-boundary}
For every formal pair $(L\subset V,W)$ with $W|_L=0$, the operations
\[
m_1=\mu_1,\qquad
m_2=m_2^{\free}+\mu_2,\qquad
m_n=\mu_n\quad (n\ge 3),
\]
define a complete filtered $A_\infty$ algebra structure on
\[
B^\intr_{L,W}:= \bigl(B^{\intr,\gr}_{L,W},\{m_n\}_{n\ge 1}\bigr).
\]
Its quasi-isomorphism class depends only on the formal germ of $W$ along $L$ and is independent of the chosen complement $V=L\oplus N$.
\end{theorem}

\begin{proof}
The only substantive issue is the Stasheff identities. Since every $\mu_n$ is a multiderivation, the full set of identities is equivalent to the vanishing of the Gerstenhaber bracket of the total Hochschild cochain
\[
\mu:=\sum_{n\ge 1}\mu_n.
\]
Under the Hochschild--Kostant--Rosenberg map, $\mu$ corresponds to the formal odd polyvector field obtained by replacing the normal variables in the Taylor expansion of $W$ by the odd derivations $\partial_{\psi_a}$. Concretely, in split coordinates
\[
\Theta_W:=\sum_{n\ge 1}
\sum_{a_1,\dots,a_n}
f_{a_1\cdots a_n}(x)\,
\partial_{\psi_{a_1}}\wedge\cdots\wedge \partial_{\psi_{a_n}}.
\]
The normal Taylor class $\Theta_W$ is a degree-$1$ element of the Gerstenhaber algebra of polyvector fields on the free graded-commutative algebra $\widehat{\Sym}(N^\vee[-1])$. In the free case $\Theta_W$ has the form $\sum_i w_i\,\partial_{\psi_i}$ where $w_i$ are formal coordinate functions on $N$ and $\partial_{\psi_i}$ are odd derivations. The Schouten--Nijenhuis bracket of two such terms is
\[
[w_i\,\partial_{\psi_i},\, w_j\,\partial_{\psi_j}]_{\sn}
= \partial_{\psi_i}(w_j)\,\partial_{\psi_j}
- \partial_{\psi_j}(w_i)\,\partial_{\psi_i},
\]
which vanishes when the $w_i$ are coordinate functions, since $\partial_{\psi_i}(w_j)=\delta_{ij}$ gives cancelling terms. At higher degree, $\Theta_W=\sum_{n\ge 1}\sum_{a_1,\ldots,a_n}f_{a_1\cdots a_n}(x)\,\partial_{\psi_{a_1}}\wedge\cdots\wedge\partial_{\psi_{a_n}}$ where the coefficients $f_{a_1\cdots a_n}(x)$ depend only on the tangent coordinates~$x$ and thus lie in the kernel of every $\partial_{\psi_a}$. The bracket $[\Theta_W,\Theta_W]_{\sn}$ reduces to a sum of terms each involving $\partial_{\psi_a}(f_{b_1\cdots b_m})=0$, so
\[
[\Theta_W,\Theta_W]_{\sn}=0.
\]
For the free completed graded-commutative algebra $B^{\intr,\gr}_{L,W}=\widehat{\Sym}(N^\vee[-1])$, the classical HKR map is strictly compatible with brackets: the HKR map is the identity on generators and extends multiplicatively, while the Gerstenhaber bracket on polyvector fields $\Lambda^*(N)$ is determined by its values on generators (degree-$1$ elements), where it coincides with the Schouten bracket. Since both brackets are biderivations determined by the same values on generators, they agree everywhere. No formality theorem is needed in the free case; the Kontsevich formality theorem is required only for non-free algebras. Transporting back through HKR gives $[\mu,\mu]_G=0$. Since $m_2^{\free}$ is associative and each $\mu_n$ is a multiderivation for $m_2^{\free}$, one also has
\[
[m_2^{\free},\mu]_G=0,
\qquad
[m_2^{\free},m_2^{\free}]_G=0.
\]
Hence the total coderivation squares to zero.

If one changes the complement $N$, the intrinsic Taylor classes $\tau_n(W)$ do not change, so the resulting $A_\infty$ structures are identified after transporting along the induced algebra isomorphism. Therefore the quasi-isomorphism class depends only on the formal germ.
\end{proof}

\begin{corollary}[Formal-neighborhood invariance; \ClaimStatusProvedHere]\label{cor:formal-neighborhood}
If two pairs $(L\subset V,W)$ and $(L'\subset V',W')$ have isomorphic formal neighborhoods carrying $W$ to $W'$, then their intrinsic boundary $A_\infty$ algebras are $A_\infty$-isomorphic.
\end{corollary}

\begin{proof}
The existence of a formal-neighborhood isomorphism is guaranteed by the formal tubular neighborhood theorem (see e.g.\ \cite[{\S}II.4]{GR14} for the derived setting): any two formal neighborhoods of a smooth submanifold in a smooth ambient space are isomorphic as filtered dg algebras, and the isomorphism is canonical up to the choice of connection. Given such an isomorphism $\varphi\colon \widehat{V}_L\xrightarrow{\sim}\widehat{V'}_{L'}$ carrying $W$ to $W'$, one obtains identifications $I_L/I_L^{n+1}\cong I_{L'}/I_{L'}^{n+1}$ at every filtration level, and therefore $\tau_n(W)=\tau_n(W')$. Formula \eqref{eq:intrinsic-mu} then produces the same higher operations, so the resulting $A_\infty$ algebras are isomorphic.
\end{proof}

\begin{corollary}[Boundary-linear strictness; \ClaimStatusProvedHere]\label{cor:boundary-linear-strict}
If all Taylor classes $\tau_n(W)$ vanish for $n\ge 2$, equivalently if $W$ is linear in the normal directions, then $B^\intr_{L,W}$ is a strict commutative dg algebra.
\end{corollary}

\begin{proof}
In that case only $\mu_1$ survives, so $m_n=0$ for $n\ge 3$ and the only deformation of the free product is the differential $m_1$.
\end{proof}

\begin{remark}
This corollary is the gateway to the exact local sector treated later. The boundary-linear Landau--Ginzburg models are exactly the models in which the intrinsic $A_\infty$ boundary algebra collapses to a strict dg algebra, making exact computation possible.
\end{remark}


\section{Multiplicativity and stabilization: Thom--Sebastiani, contractible blocks, and Clifford blocks}\label{sec:ts-knorrer}

The intrinsic boundary construction should behave like a stable algebraic object. Three operations test this stability: external sums, tangent-normal hyperbolic stabilizations, and pure-normal hyperbolic stabilizations.

\subsection{Thom--Sebastiani}

Let $(L_i\subset V_i,W_i)$, $i=1,2$, be two formal pairs with $W_i|_{L_i}=0$. Form
\[
(L_1\oplus L_2 \subset V_1\oplus V_2,\; W_1\boxplus W_2),
\qquad
(W_1\boxplus W_2)(v_1,v_2)=W_1(v_1)+W_2(v_2).
\]

\begin{theorem}[Boundary Thom--Sebastiani; \ClaimStatusProvedHere]\label{thm:TS}
There is a natural filtered $A_\infty$ isomorphism
\[
B^\intr_{L_1\oplus L_2,\;W_1\boxplus W_2}
\cong
B^\intr_{L_1,W_1}\,\widehat{\otimes}\, B^\intr_{L_2,W_2}.
\]
\end{theorem}

\begin{proof}
The product neighborhood $N(L_1\times L_2)\subset M_1\times M_2$ splits as $N(L_1)\times N(L_2)$ because the normal bundle of a product embedding is the external direct sum $N_{L_1\times L_2/M_1\times M_2}\cong N_{L_1/M_1}\boxplus N_{L_2/M_2}$. The ideal filtration therefore decomposes as
\[
I^n(L_1\times L_2)=\sum_{p+q=n} I^p(L_1)\otimes I^q(L_2)
\]
by the K\"unneth formula for completed tensor products. The normal Taylor class is additive:
\[
\Theta_{W_1\boxplus W_2}=\Theta_{W_1}\otimes 1+1\otimes \Theta_{W_2},
\]
because the Taylor expansion of the external-sum potential $W_1(x)+W_2(y)$ in normal coordinates separates into independent summands. Consequently the individual Taylor classes satisfy $\tau_n(W_1\boxplus W_2)=\tau_n(W_1)\oplus \tau_n(W_2)$, and the multiderivations defining the $A_\infty$ structure on the left are exactly the sums of the multiderivations from the two factors, which is the completed tensor product $A_\infty$ structure.
\end{proof}

\subsection{Tangent-normal hyperbolic stabilization}

Suppose $U$ is a finite-dimensional vector space and consider the enlarged data
\[
V':=V\oplus U\oplus U^\vee,
\qquad
L':=L\oplus U,
\qquad
W':=W+\langle u,v\rangle,
\]
where $u\in U$ is tangent to the boundary and $v\in U^\vee$ is normal. Then one leg of the hyperbolic pair is tangent and one leg is normal.

\begin{theorem}[Contractible stabilization; \ClaimStatusProvedHere]\label{thm:contractible-stab}
With the above notation,
\[
B^\intr_{L',W'}
\cong
B^\intr_{L,W}\,\widehat{\otimes}\,
\bigl(\widehat{\Sym}(U^\vee)\ot \Lambda(U^\vee),\, d\xi=u\bigr),
\]
where the second factor is the ordinary completed Koszul resolution of $\kk$. In particular
\[
B^\intr_{L',W'}\simeq B^\intr_{L,W}
\]
as $A_\infty$ algebras.
\end{theorem}

\begin{proof}
The new potential is the external sum of $W$ and the linear-normal function $\langle u,v\rangle$ on $U\oplus U^\vee$ relative to the boundary $U\subset U\oplus U^\vee$. By Theorem~\ref{thm:TS} it suffices to compute the second factor. But that factor is exactly the strict dg algebra with even variables $u$ and odd variables $\xi$ corresponding to the normal directions, with differential $d\xi=u$. This is contractible.
\end{proof}

\subsection{Pure-normal hyperbolic stabilization}

Now suppose both legs of the hyperbolic pair are normal:
\[
V' := V\oplus U\oplus U^\vee,
\qquad
L' := L,
\qquad
W' := W+\langle u,v\rangle,
\]
where both $u$ and $v$ are normal directions.

\begin{theorem}[Hyperbolic Clifford stabilization; \ClaimStatusProvedHere]\label{thm:clifford-stab}
There is a natural dg isomorphism
\[
B^\intr_{L,W+\langle u,v\rangle}
\cong
B^\intr_{L,W}\,\widehat{\otimes}\, \Cl(U\oplus U^\vee),
\]
where $\Cl(U\oplus U^\vee)$ denotes the completed hyperbolic Clifford algebra. Consequently
\[
\Perf\bigl(B^\intr_{L,W+\langle u,v\rangle}\bigr)
\simeq
\Perf\bigl(B^\intr_{L,W}\bigr)
\]
after stabilization by the spinor module $\Lambda(U)$.
\end{theorem}

\begin{proof}
In the pure-normal case, the new Taylor data are quadratic and nondegenerate entirely in the normal variables. The intrinsic $A_\infty$ operations contributed by the pair $(u,v)$ assemble into the standard hyperbolic Clifford relation. The resulting factor is exactly $\Cl(U\oplus U^\vee)$. Since this Clifford algebra is Morita equivalent to $\kk$ via the spinor module $\Lambda(U)$, tensoring with it does not change the perfect module category.
\end{proof}

\begin{remark}
This is the boundary-level zero-mode shadow of Kn\"orrer periodicity. What matters is not a periodicity statement in a triangulated category, but the explicit identification of the open-sector factor that is being added: a hyperbolic Clifford block.
\end{remark}

\subsection{Relative formal Morse splitting}

The preceding theorems combine with a formal Morse lemma relative to the boundary.

\begin{theorem}[Relative formal Morse splitting; \ClaimStatusProvedHere]\label{thm:relative-morse}
Let $V=L\oplus K\oplus U\oplus U^\vee$ and let
\[
W(x,k,u,v)
\]
be a formal potential vanishing on $L$ whose quadratic part in the $(u,v)$ variables is the split hyperbolic form $\langle u,v\rangle$. Assume there are no linear terms in $u$ or $v$ and no mixed quadratic terms between $(u,v)$ and the residual normal directions $K$. Then there is a formal automorphism preserving $L$ and fixing the residual normal directions to first order such that
\[
W \sim W_{\mathrm{red}}(x,k)+\langle u,v\rangle.
\]
Consequently
\[
B^\intr_{L,W}
\simeq
B^\intr_{L,W_{\mathrm{red}}}\,\widehat{\otimes}\,\Cl(U\oplus U^\vee)
\]
in the pure-normal case, and analogously with a contractible Koszul factor in the tangent-normal case.
\end{theorem}

\begin{proof}
The formal automorphism is constructed inductively by degree. Because the quadratic form on $(u,v)$ is split and nondegenerate, the homological operator given by contraction with $\langle u,v\rangle$ is invertible on all higher monomials involving at least one $u$ or $v$. Therefore every higher interaction involving the hyperbolic directions can be killed order by order. The factorization of the boundary algebra then follows from Theorems~\ref{thm:contractible-stab} and \ref{thm:clifford-stab}.
\end{proof}

\begin{remark}
Split hyperbolic directions are explicitly removable; what remains is the genuine singular residue. This reduction drives the exact sector below.
\end{remark}


\section{The exact boundary-linear Landau--Ginzburg sector}\label{sec:boundary-linear}

The passage to the exact sector where the intrinsic $A_\infty$ algebra becomes a strict dg algebra and every object can be written explicitly.

Let $V=L\oplus N$ and let
\[
F:L\longrightarrow N^\vee
\]
be a formal map. Define the superpotential
\[
W(x,y):=\langle y, F(x)\rangle,
\qquad
x\in L,\; y\in N.
\]
Then $W$ vanishes on $L$ and is linear in the normal variables $y$.

\subsection{Boundary dg algebra and derived zero locus}

Choose a basis $n_a$ of $N$ and odd generators $\eta_a$ dual to $n_a$. The boundary algebra is the strict commutative dg algebra
\begin{equation}\label{eq:boundary-linear-dga}
B_{L,W}
:=
\bigl(\widehat{\cO}(L)\ot \Lambda(N),\ d\eta_a=F_a(x),\ dx_i=0\bigr).
\end{equation}
This is the Koszul algebra of the equations $F_a(x)=0$.

\begin{definition}
The derived boundary zero locus of $F$ is
\[
Z_F^{\der}:=\Spec(B_{L,W}).
\]
\end{definition}

\begin{remark}
In ordinary geometry $Z_F$ would be the vanishing locus $F^{-1}(0)\subset L$. The dg algebra $B_{L,W}$ remembers the derived structure and therefore the higher obstruction theory of the equations.
\end{remark}

\subsection{The bulk critical locus as a shifted cotangent}

The bulk derived critical locus of $W$ is the dg scheme cut out by the equations
\[
\frac{\partial W}{\partial y_a}=F_a(x)=0,
\qquad
\frac{\partial W}{\partial x_i}
=
\sum_a y_a\frac{\partial F_a}{\partial x_i}(x)=0.
\]
This is the shifted cotangent to the derived zero locus.

\begin{theorem}[Derived critical locus theorem; \ClaimStatusProvedHere]\label{thm:dcrit-shiftedcot}
For the boundary-linear superpotential $W(x,y)=\langle y,F(x)\rangle$ there is a canonical equivalence of formal dg schemes
\[
\dCrit(W)\simeq T^*[-1] Z_F^{\der}.
\]
\end{theorem}

\begin{proof}
The shifted cotangent algebra of $Z_F^{\der}$ is the algebra of polyvector fields on $Z_F^{\der}$:
\[
\cO\bigl(T^*[-1]Z_F^{\der}\bigr)\cong \PV(Z_F^{\der}).
\]
Using the Koszul model \eqref{eq:boundary-linear-dga}, this algebra is generated by:
\begin{itemize}[leftmargin=1.8em]
\item even coordinates $x_i$ on $L$,
\item odd coordinates $\eta_a$ from the derived equations $F_a(x)=0$,
\item odd tangent polyvectors $\theta_i$ dual to $x_i$,
\item even tangent polyvectors $y_a$ dual to $\eta_a$.
\end{itemize}
Its differential is induced by the boundary differential and is exactly
\[
d\theta_i
=
\sum_a y_a\frac{\partial F_a}{\partial x_i}(x),
\qquad
dy_a=0,
\qquad
d\eta_a=F_a(x),
\qquad
dx_i=0.
\]
But this is exactly the standard dg algebra of functions on the derived critical locus of $W(x,y)=\langle y,F(x)\rangle$. Hence the two dg schemes are canonically equivalent.
\end{proof}

\subsection{Derived center of the boundary algebra}

The usual HKR theorem identifies Hochschild cochains of a smooth commutative algebra with polyvector fields. In the present dg setting, the same philosophy gives the derived center of the boundary dg algebra as the shifted cotangent algebra of the derived boundary zero locus. Because the bulk critical locus is exactly that shifted cotangent, the bulk is recovered from the boundary algebra.

\begin{theorem}[Boundary-linear bulk/boundary theorem; \ClaimStatusProvedHere]\label{thm:boundary-linear-bulk-boundary}
For $W(x,y)=\langle y,F(x)\rangle$ there is a canonical quasi-isomorphism
\[
\Zder(B_{L,W})
\simeq
\HH^\bullet(B_{L,W})
\simeq
\cO\bigl(T^*[-1]Z_F^{\der}\bigr)
\simeq
\cO\bigl(\dCrit(W)\bigr).
\]
In particular the derived center of the boundary algebra equals the bulk algebra in the boundary-linear sector.
\end{theorem}

\begin{proof}
Because $B_{L,W}$ is a cofibrant completed commutative dg algebra, its derived center is modeled by Hochschild cochains. The algebra $B_{L,W} = (\widehat{\cO}(L) \otimes \Lambda(N), d\eta_a = F_a(x))$ is quasi-smooth: its cotangent complex $\mathbb{L}_{B/\kk}$ is concentrated in degrees $[0,1]$, where the degree-$0$ part is $\Omega^1_{B/\kk}$ (K\"ahler differentials of the underlying commutative algebra) and the degree-$1$ part is the normal bundle $N$. This two-term structure arises because $L$ is a smooth submanifold of the smooth ambient space $V = L \oplus N$, so the conormal sequence $0 \to N^\vee \to \Omega^1_V|_L \to \Omega^1_L \to 0$ yields a two-term cotangent complex; since both $L$ and $N$ are finitely generated free modules over the formal power series ring, this complex is perfect. The derived HKR isomorphism therefore applies (see Calaque--Van den Bergh \cite{CVdB10} for the dg setting, or \cite[Theorem~7.8.3.3]{Lur18} for the derived algebraic geometry framework), identifying Hochschild cochains with polyvector fields on the derived affine dg scheme $Z_F^{\der}$. Theorem~\ref{thm:dcrit-shiftedcot} identifies those polyvectors with functions on the bulk critical locus.
\end{proof}

\subsection{The free polarization theorem as the zero superpotential case}

Setting $F=0$ gives the free model. Then
\[
B_L^{\free}
=
\widehat{\Sym}(L^\vee)\ot \Lambda(N),
\qquad
Z_0^{\der}=L\oplus N^\vee[1].
\]
The bulk algebra is
\[
A_{\mathrm{bulk}}^{\free}(V)
=
\widehat{\Sym}(V^\vee)\ot \Lambda(V).
\]

\begin{corollary}[Free polarization theorem; \ClaimStatusProvedHere]\label{cor:free-polarization}
For a linear boundary condition $L\subset V$ with quotient $Q=V/L$ one has
\[
B_L^{\free}\cong \widehat{\Sym}(L^\vee)\ot \Lambda(Q),
\qquad
K_L^{\free}\simeq \widehat{\Sym}(Q^\vee)\ot \Lambda(L),
\]
and
\[
\HH^\bullet(B_L^{\free})
\cong
\HH^\bullet(K_L^{\free})
\cong
\widehat{\Sym}(V^\vee)\ot \Lambda(V).
\]
Thus the boundary and line algebras are opposite Koszul polarizations of the full free bulk algebra.
\end{corollary}

\begin{proof}
The first identity is the boundary dg algebra with $F=0$. The second is ordinary completed Koszul duality. The final statement follows from Theorem~\ref{thm:boundary-linear-bulk-boundary} and the Morita invariance of Hochschild cohomology.
\end{proof}


\section{Exact pointed lines, one-step Jacobi coalgebras, and vacuum twisting}\label{sec:exact-lines}

The boundary-linear sector is exact enough not only to compute the boundary algebra but also to compute the line algebra attached to a \emph{pointed} boundary vacuum.

Fix a point
\[
p\in F^{-1}(0)\subset L.
\]
Complete at $p$ and set
\[
U:=T_pL,
\qquad
R:=N^\vee.
\]
The completed local boundary algebra is
\[
B_{F,p}
=
\bigl(\kk[[U^\vee]]\ot \Lambda(R^\vee),\ d\eta_\alpha=F_\alpha(x)\bigr).
\]

\subsection{The one-step Jacobi coalgebra}

The Taylor coefficients of $F$ at $p$ assemble into a canonical dg coalgebra.

\begin{definition}[One-step Jacobi coalgebra]
The one-step Jacobi coalgebra of the pointed germ $(F,p)$ is
\[
\cJ_{F,p}:=\bigl(S^c(sU\oplus R),\, b_F\bigr),
\]
where $S^c$ denotes the cofree cocommutative coalgebra and the coderivation $b_F$ is determined by its Taylor coefficients
\[
b_{F,n}: S^n(sU)\longrightarrow R
\]
given by the $n$-th derivatives of $F$ at $p$.
\end{definition}

The terminology is deliberate: this coalgebra is the local coalgebraic envelope of the boundary equations. It is the object from which both the local boundary algebra and the pointed line algebra can be recovered.

\begin{theorem}[Exact presentation theorem; \ClaimStatusProvedHere]\label{thm:exact-line}
The completed boundary algebra is canonically dual to the one-step
Jacobi coalgebra, and the exact pointed line algebra is its cobar
algebra. Equivalently, the pointed line algebra computes the derived
endomorphism algebra of the residue field over the completed local
boundary algebra.
\end{theorem}

\begin{proof}
The completed dual of the cofree coalgebra $S^c(sU\oplus R)$ is the completed free commutative algebra on $U^\vee$ and $R^\vee[-1]$, which is exactly the underlying graded algebra of $B_{F,p}$. The coderivation $b_F$ dualizes to the differential $d\eta_\alpha=F_\alpha(x)$, so $B_{F,p}\cong \cJ_{F,p}^\vee$.

Standard bar--cobar/Koszul duality for complete augmented commutative dg algebras identifies the derived endomorphisms of the residue field with the cobar of the dual coalgebra. Applying this to $B_{F,p}$ yields
\[
\RHom_{B_{F,p}}(\kk,\kk)\simeq \Cobar(\cJ_{F,p}).
\]
\end{proof}

\subsection{Explicit formula for the pointed line algebra}

Choose coordinates $x_1,\dots,x_u$ on $U$ centered at $p$, a basis $r_\alpha$ of $R$, odd generators $\lambda_i$ of degree $1$ dual to $x_i$, and even generators $c_\alpha$ of degree $2$ dual to $r_\alpha$.

Define the graded algebra
\[
K_{F,p}^{\gr}:=\Sym(R)\ot \Lambda(U)=\kk[c_\alpha]\ot \Lambda(\lambda_i).
\]
For $n\ge 1$ define Hochschild cochains
\begin{equation}\label{eq:mu-F-explicit}
\mu_n(a_1,\dots,a_n)
:=
\sum_{\alpha}
c_\alpha
\sum_{i_1,\dots,i_n}
\frac{1}{n!}
\frac{\partial^n F_\alpha}{\partial x_{i_1}\cdots \partial x_{i_n}}(p)
\bigl(\partial_{\lambda_{i_1}}a_1\bigr)\cdots \bigl(\partial_{\lambda_{i_n}}a_n\bigr).
\end{equation}
Set
\[
m_1=\mu_1,
\qquad
m_2=m_2^{\free}+\mu_2,
\qquad
m_n=\mu_n\quad (n\ge 3).
\]

\begin{corollary}[Explicit pointed line algebra; \ClaimStatusProvedHere]\label{cor:explicit-line}
The operations above define a complete filtered $A_\infty$ algebra
structure on $K_{F,p}^{\gr}$, denoted $K^{\expmod}_{F,p}$. This
algebra is $A_\infty$ quasi-isomorphic to the exact pointed line algebra
and to the derived endomorphism algebra of the residue field over the
completed boundary algebra.
\end{corollary}

\begin{proof}
Under the HKR map, the total cochain $\mu=\sum_n \mu_n$ corresponds to the odd polyvector field obtained from the Taylor expansion of $F$ by replacing the tangent variables with the odd derivations $\partial_{\lambda_i}$. Since these odd derivations supercommute and the coefficients are scalar functions, the Schouten bracket vanishes. Exactly as in the proof of Theorem~\ref{thm:intrinsic-boundary}, this implies the Stasheff identities. The comparison with $\Cobar(\cJ_{F,p})$ is tautological from the construction of $b_F$.
\end{proof}

\subsection{Vacua as Maurer--Cartan elements}

One of the most satisfying features of the one-step Jacobi coalgebra is that classical vacua become internal Maurer--Cartan data.

\begin{theorem}[Maurer--Cartan vacuum theorem; \ClaimStatusProvedHere]\label{thm:mc-vacua}
Let $q$ be a formal point of $L$ near $p$. Then the following are equivalent:
\begin{enumerate}[label=(\roman*),leftmargin=1.8em]
\item $q$ is a boundary vacuum, that is, $F(q)=0$;
\item the displacement $u=q-p\in U$ defines a Maurer--Cartan element of the coalgebra $\cJ_{F,p}$;
\item $u$ defines a Maurer--Cartan element of the exact pointed line algebra.
\end{enumerate}
Twisting $\cJ_{F,p}$ by the Maurer--Cartan element corresponding to $q$ yields the recentered coalgebra:
\[
\cJ_{F,p}^{\,q}\cong \cJ_{F,q}.
\]
\end{theorem}

\begin{proof}
The Maurer--Cartan equation in $\cJ_{F,p}$ is
\[
\sum_{n\ge 1}\frac{1}{n!} b_{F,n}(u,\dots,u)=0,
\]
which is precisely the Taylor expansion of $F(p+u)=0$. The equivalence with the $A_\infty$ Maurer--Cartan equation follows under the bar--cobar correspondence. Twisting by $u$ simply translates the Taylor expansion from the base point $p$ to the base point $q$.
\end{proof}

\subsection{Polarized reconstruction}

The exact line algebra remembers more than its quasi-isomorphism class: it remembers a primitive polarization.

\begin{definition}[Primitive polarization]
The primitive part of the bar coalgebra of the exact pointed line algebra
is canonically $sU\oplus R$.
A \emph{polarization} is the decomposition of this primitive space into its odd tangent piece $sU$ and its even obstruction piece $R$.
\end{definition}

\begin{theorem}[Polarized reconstruction; \ClaimStatusProvedHere]\label{thm:polarized-reconstruction}
The polarized bar coalgebra of the exact pointed line algebra reconstructs
the full pointed formal germ of the map $F$ at $p$. Equivalently, the
polarized coalgebra $\cJ_{F,p}$ is a complete invariant of the pointed
boundary-linear formal germ.
\end{theorem}

\begin{proof}
The coderivation of $\cJ_{F,p}$ is determined exactly by the maps
\[
b_{F,n}:S^n(sU)\to R.
\]
These are the Taylor coefficients of $F$ at $p$. Conversely, given the polarized coalgebra, the maps $b_{F,n}$ may be read off from the component of the coderivation landing in the $R$-summand of the primitive part. Reassembling these Taylor coefficients recovers the formal map $F$.
\end{proof}

\begin{remark}
This theorem is the first local incarnation of the guiding slogan that the line sector should not merely receive data from the boundary equations; it should \emph{present} them.
\end{remark}


\section{Boundary Kuranishi elimination and the minimal pointed line algebra}\label{sec:kuranishi}

The exact pointed line algebra is canonical, but it is not yet minimal. If the linearization of $F$ at $p$ is nonzero, then the line algebra has a nontrivial unary operation. The next theorem integrates out the linearly massive directions exactly and produces a minimal local model.

\subsection{Splitting the linearized equations}

Let
\[
dF_p:U\to R.
\]
Choose splittings
\[
U=T\oplus M,
\qquad
R=I\oplus C,
\]
with
\[
T=\ker(dF_p),
\qquad
I=\im(dF_p),
\qquad
C\cong \coker(dF_p),
\]
and such that
\[
A:=dF_p|_M:M\xrightarrow{\sim} I
\]
is an isomorphism.

Choose coordinates $(u,v)$ on $T\oplus M$ and decompose
\[
F=(F_I,F_C),
\qquad
F_I\in I[[u,v]],
\quad
F_C\in C[[u,v]].
\]

\begin{lemma}[Formal implicit solution; \ClaimStatusProvedHere]\label{lem:formal-implicit}
There exists a unique formal map
\[
\varphi(u)\in M[[u]]
\]
with no constant or linear term such that
\[
F_I\bigl(u,\varphi(u)\bigr)=0.
\]
\end{lemma}

\begin{proof}
Expand $\varphi(u)=\sum_{n\ge 2}\varphi_n(u)$ by total degree in $u$. At each order $n$, the equation takes the form
\[
A\,\varphi_n + (\text{known term of degree }n)=0.
\]
Since $A$ is invertible, $\varphi_n$ is uniquely determined recursively.
\end{proof}

\begin{definition}[Effective obstruction map]
The reduced Kuranishi map is
\begin{equation}\label{eq:kappa-def}
\kappa(u):=F_C\bigl(u,\varphi(u)\bigr)\in C[[u]].
\end{equation}
It has no constant or linear term.
\end{definition}

\subsection{Exact splitting of the boundary dg algebra}

Set
\[
w:=v-\varphi(u).
\]
Then $F_I(u,\varphi(u)+w)$ vanishes at $w=0$, so it factors by $w$.

\begin{lemma}[Factorization of the massive equations; \ClaimStatusProvedHere]\label{lem:factorization}
There exist formal matrix-valued power series
\[
\mathcal A(u,w)\in \End(M,I)[[u,w]],
\qquad
\mathcal B(u,w)\in \Hom(M,C)[[u,w]],
\]
with $\mathcal A(0,0)=A$ invertible, such that
\[
F_I\bigl(u,\varphi(u)+w\bigr)=\mathcal A(u,w)\,w,
\qquad
F_C\bigl(u,\varphi(u)+w\bigr)=\kappa(u)+\mathcal B(u,w)\,w.
\]
\end{lemma}

\begin{proof}
This is the formal Taylor expansion in the $w$-variables around $w=0$, using the fact that the constant term of $F_I$ vanishes by definition of $\varphi$.
\end{proof}

\begin{theorem}[Boundary Kuranishi elimination; \ClaimStatusProvedHere]\label{thm:boundary-kuranishi}
After the changes of variables
\[
v=\varphi(u)+w
\]
and the odd change of generators
\[
\xi:=\mathcal A(u,w)^{-1}\eta_I,
\qquad
\zeta:=\eta_C-\mathcal B(u,w)\mathcal A(u,w)^{-1}\eta_I,
\]
the completed local boundary dg algebra becomes
\[
B_{F,p}
\cong
\bigl(\kk[[u,w]]\ot \Lambda(\xi,\zeta),\ d\xi=w,\ d\zeta=\kappa(u)\bigr).
\]
Consequently there is a canonical dg quasi-isomorphism
\[
B_{F,p}\simeq B_\kappa
:=
\bigl(\kk[[u]]\ot \Lambda(C^\vee),\ d\zeta=\kappa(u)\bigr).
\]
\end{theorem}

\begin{proof}
The displayed formula follows by a direct calculation:
\[
d\xi=\mathcal A^{-1}d\eta_I=\mathcal A^{-1}\mathcal A w=w,
\]
and
\[
d\zeta=d\eta_C-\mathcal B\mathcal A^{-1}d\eta_I
=(\kappa+\mathcal Bw)-\mathcal B\mathcal A^{-1}\mathcal A w
=\kappa.
\]
The factor $\bigl(\kk[[w]]\ot \Lambda(\xi),d\xi=w\bigr)$ is the completed Koszul resolution of $\kk$ and is therefore contractible. Removing it gives the reduced dg algebra $B_\kappa$.
\end{proof}

\begin{remark}
This theorem is stronger than ordinary homological perturbation. It does not merely assert that a smaller model exists. It constructs an explicit nonlinear change of generators that exhibits the massive directions as a contractible factor. For that reason it is genuinely a nonlinear Schur complement.
\end{remark}

\subsection{The minimal pointed line algebra}

Apply the exact line construction to the reduced map $\kappa:T\to C$.

Let
\[
K_\kappa^{\gr}:=\Sym(C)\ot \Lambda(T).
\]
For $n\ge 2$ define
\begin{equation}\label{eq:mu-kappa}
\mu_n^\kappa(a_1,\dots,a_n)
:=
\sum_{\alpha}
c_\alpha
\sum_{i_1,\dots,i_n}
\frac{1}{n!}
\frac{\partial^n\kappa_\alpha}{\partial u_{i_1}\cdots \partial u_{i_n}}(0)
\bigl(\partial_{\lambda_{i_1}}a_1\bigr)\cdots \bigl(\partial_{\lambda_{i_n}}a_n\bigr),
\end{equation}
and set
\[
m_1^\kappa=0,
\qquad
m_2^\kappa=m_2^{\free}+\mu_2^\kappa,
\qquad
m_n^\kappa=\mu_n^\kappa\quad (n\ge 3).
\]

\begin{theorem}[Minimal pointed line algebra; \ClaimStatusProvedHere]\label{thm:minimal-pointed-line}
The operations above define a minimal $A_\infty$ algebra
\[
K_\kappa
=
\bigl(\Sym(C)\ot \Lambda(T),\{m_n^\kappa\}_{n\ge 2}\bigr),
\]
and there are $A_\infty$ quasi-isomorphisms
\[
K_{\mathrm{ex},F,p}
\simeq
K_\kappa.
\]
Equivalently, the exact pointed line algebra is governed, up to quasi-isomorphism, by the reduced Kuranishi map $\kappa$.
\end{theorem}

\begin{proof}
Theorem~\ref{thm:boundary-kuranishi} identifies $B_{F,p}$ with the tensor product of $B_\kappa$ and a contractible Koszul factor. Derived endomorphisms of the residue field are invariant under tensoring with such a factor, so
\[
\RHom_{B_{F,p}}(\kk,\kk)\simeq \RHom_{B_\kappa}(\kk,\kk).
\]
Applying Theorem~\ref{thm:exact-line} to $\kappa$ identifies the right-hand side with the explicit $A_\infty$ algebra above. Since $\kappa$ has no linear term, the resulting $A_\infty$ algebra is minimal.
\end{proof}

\subsection{Recursive formulas and the first effective couplings}

The reduced map $\kappa$ admits exact recursive formulas. Write
\[
\Phi(u):=(u,\varphi(u))=\sum_{n\ge 1}\frac{1}{n!}\Phi_n(u,\dots,u),
\qquad
\kappa(u)=\sum_{n\ge 2}\frac{1}{n!}\kappa_n(u,\dots,u).
\]
Write the Taylor expansion of $F_I$ and $F_C$ at the origin as
\[
F_I=\sum_{q\ge 1}F_q^I,
\qquad
F_C=\sum_{q\ge 1}F_q^C.
\]

\begin{proposition}[Recursive Kuranishi coefficients; \ClaimStatusProvedHere]\label{prop:recursive-kuranishi}
The multilinear maps $\Phi_n$ and $\kappa_n$ are uniquely determined by
\begin{align}
A\Phi_n
&=
-
\sum_{q=2}^{n}
\frac{1}{q!}
\sum_{\substack{n_1+\cdots+n_q=n\\ n_i\ge 1}}
\operatorname{Sym}\Bigl(F_q^I\bigl(\Phi_{n_1},\dots,\Phi_{n_q}\bigr)\Bigr),
\label{eq:phi-recursion}
\\[0.5em]
\kappa_n
&=
\sum_{q=2}^{n}
\frac{1}{q!}
\sum_{\substack{n_1+\cdots+n_q=n\\ n_i\ge 1}}
\operatorname{Sym}\Bigl(F_q^C\bigl(\Phi_{n_1},\dots,\Phi_{n_q}\bigr)\Bigr).
\label{eq:kappa-recursion}
\end{align}
\end{proposition}

\begin{proof}
Insert the exponential expansions into the defining equations
\[
F_I(\Phi(u))=0,
\qquad
\kappa(u)=F_C(\Phi(u)),
\]
and compare homogeneous components of degree $n$.
\end{proof}

\begin{corollary}[Quadratic and cubic effective couplings; \ClaimStatusProvedHere]\label{cor:low-order-kappa}
The first two Kuranishi coefficients are
\begin{align}
\kappa_2(t_1,t_2)
&=
F_2^C(t_1,t_2),
\label{eq:kappa2}
\\[0.5em]
\kappa_3(t_1,t_2,t_3)
&=
F_3^C(t_1,t_2,t_3)
-
\sum_{\mathrm{cyc}}
F_2^C\Bigl(t_1,\,A^{-1}F_2^I(t_2,t_3)\Bigr).
\label{eq:kappa3}
\end{align}
Consequently
\[
m_2^\kappa(\lambda_{t_1},\lambda_{t_2})
=
\lambda_{t_1}\wedge \lambda_{t_2}
+\frac{1}{2}\kappa_2(t_1,t_2),
\qquad
m_3^\kappa(\lambda_{t_1},\lambda_{t_2},\lambda_{t_3})
=
\frac{1}{6}\kappa_3(t_1,t_2,t_3).
\]
\end{corollary}

\begin{proof}
At quadratic order only the quadratic target component $F_2^C$ contributes. At cubic order there are two kinds of contributions: the direct cubic term $F_3^C$ and the indirect term obtained by feeding two tangent inputs into the massive channel through $F_2^I$, propagating back via $A^{-1}$, and returning to the obstruction space through $F_2^C$.
\end{proof}

\begin{remark}
Formula \eqref{eq:kappa3} is the algebraic avatar of integrating out a massive propagator. The effective cubic interaction is not the raw cubic Taylor coefficient of the original equations; it is the raw cubic term corrected by the tree-level exchange through the massive sector.
\end{remark}

\subsection{Smooth vacua}

\begin{corollary}[Smooth-point formality; \ClaimStatusProvedHere]\label{cor:smooth-formality}
If $p$ is a smooth point of the classical zero locus $F^{-1}(0)$, then $C=0$ and the minimal pointed line algebra is the exterior algebra
\[
K_\kappa\simeq \Lambda(T_pF^{-1}(0)).
\]
\end{corollary}

\begin{proof}
Smoothness is equivalent to surjectivity of $dF_p$, hence $C=\coker(dF_p)=0$. Therefore $\kappa=0$ and only the free exterior algebra on $T=\ker(dF_p)$ remains.
\end{proof}


\section{Local bulk from local lines, and the boundary-linear line Mather--Yau theorem}\label{sec:classification}

The previous section produces the minimal pointed line algebra $K_\kappa$ attached to a boundary vacuum. The proof of two culmination theorems:
\begin{enumerate}[label=(\arabic*),leftmargin=1.8em]
\item the local bulk algebra is recovered from $K_\kappa$ by Hochschild cochains;
\item the polarized quasi-isomorphism class of $K_\kappa$ classifies the reduced microlocal formal germ.
\end{enumerate}

\subsection{Local bulk from local lines}

Let $\gamma$ denote the coordinate on $C^\vee$. The effective boundary-linear potential determined by the Kuranishi map is
\[
W_{\eff}(u,\gamma):=\langle \gamma,\kappa(u)\rangle.
\]
Its boundary dg algebra is exactly $B_\kappa$.

\begin{theorem}[Local bulk/line theorem; \ClaimStatusProvedHere]\label{thm:local-bulk-line}
For every pointed boundary-linear germ $(F,p)$ one has canonical quasi-isomorphisms
\[
\HH^\bullet(K_{F,p})
\simeq
\HH^\bullet(K_\kappa)
\simeq
\HH^\bullet(B_\kappa)
\simeq
\cO\bigl(\dCrit(W_{\eff})\bigr).
\]
Equivalently, the completed local bulk algebra depends only on the minimal pointed line algebra.
\end{theorem}

\begin{proof}
The first equivalence is Theorem~\ref{thm:minimal-pointed-line}. Since $K_\kappa$ and $B_\kappa$ are Koszul dual presentations of the same pointed local category, their Hochschild cochains agree by Morita invariance. The last equivalence is the boundary-linear bulk/boundary theorem applied to the reduced map $\kappa:T\to C$.
\end{proof}

\begin{remark}
This is the exact local shadow of the global triangle \eqref{eq:global-corrected-triangle}. The line algebra does not equal the bulk algebra; the bulk is the derived center of the line algebra, and in the exact local sector that center is computable.
\end{remark}

\subsection{Polarized quasi-isomorphisms}

The reconstruction theorem showed that the polarized bar coalgebra of the exact pointed line algebra recovers the pointed formal germ. After Kuranishi reduction the same statement becomes stronger: it classifies the \emph{reduced} microlocal germ, which is exactly the singular residue left after all linearly massive directions have been eliminated.

\begin{definition}[Polarized $A_\infty$ isomorphism]
Let
\[
K_\kappa=\Sym(C)\ot \Lambda(T),
\qquad
K_{\kappa'}=\Sym(C')\ot \Lambda(T')
\]
be minimal pointed line algebras. A \emph{polarized $A_\infty$ isomorphism}
\[
\Phi:K_\kappa\longrightarrow K_{\kappa'}
\]
is an $A_\infty$ isomorphism whose induced bar-coalgebra isomorphism
\[
\BBar \Phi:\BBar K_\kappa\to \BBar K_{\kappa'}
\]
preserves the primitive decomposition
\[
\mathrm{Prim}(\BBar K_\kappa)=sT\oplus C,
\qquad
\mathrm{Prim}(\BBar K_{\kappa'})=sT'\oplus C'.
\]
\end{definition}

\begin{definition}[Formal equivalence of reduced germs]
Two reduced Kuranishi maps
\[
\kappa:T\to C,
\qquad
\kappa':T'\to C'
\]
are \emph{formally equivalent} if there exist formal linear isomorphisms on source and target, together with higher formal changes of tangent coordinates preserving the base point, carrying the Taylor expansion of $\kappa$ to that of $\kappa'$.
\end{definition}

\begin{theorem}[Boundary-linear line Mather--Yau theorem; \ClaimStatusProvedHere]\label{thm:line-mather-yau}
Let $(F,p)$ and $(G,q)$ be pointed boundary-linear germs with reduced Kuranishi maps
\[
\kappa_F:T_F\to C_F,
\qquad
\kappa_G:T_G\to C_G.
\]
Then the following are equivalent:
\begin{enumerate}[label=(\roman*),leftmargin=1.8em]
\item $\kappa_F$ and $\kappa_G$ are formally equivalent;
\item the minimal pointed line algebras $K_{\kappa_F}$ and $K_{\kappa_G}$ are polarized $A_\infty$-isomorphic;
\item the one-step Jacobi coalgebras $\cJ_{\kappa_F}$ and $\cJ_{\kappa_G}$ are polarized isomorphic.
\end{enumerate}
Whenever these conditions hold, the completed local bulk algebras are canonically quasi-isomorphic:
\[
\cO\bigl(\dCrit(W_{\eff,F})\bigr)\simeq \cO\bigl(\dCrit(W_{\eff,G})\bigr).
\]
\end{theorem}

\begin{proof}
\emph{(i)$\Rightarrow$(iii).}
A formal equivalence of reduced germs identifies all Taylor coefficients of $\kappa_F$ with those of $\kappa_G$ after transport by the induced source and target changes. By definition these Taylor coefficients are exactly the structure maps of the one-step Jacobi coalgebras, so the coalgebras are polarized isomorphic.

\emph{(iii)$\Rightarrow$(ii).}
Cobar sends a polarized coalgebra isomorphism to an $A_\infty$ isomorphism between the corresponding cobar algebras. But $K_{\kappa_F}\simeq \Cobar(\cJ_{\kappa_F})$ and similarly for $G$.

\emph{(ii)$\Rightarrow$(iii).}
Apply the bar construction. Because the algebras are minimal, the bar construction recovers a complete filtered coalgebra with primitive part equal to $sT\oplus C$. The polarization hypothesis guarantees that the induced isomorphism preserves the tangent and obstruction summands separately, so it is a polarized coalgebra isomorphism.

\emph{(iii)$\Rightarrow$(i).}
By the polarized reconstruction theorem, a polarized coalgebra recovers the full reduced formal map. Therefore a polarized coalgebra isomorphism yields a formal equivalence of Kuranishi maps.

The final statement about bulk algebras follows from Theorem~\ref{thm:local-bulk-line}.
\end{proof}

\begin{corollary}[Microlocal classification theorem; \ClaimStatusProvedHere]\label{cor:microlocal-classification}
The assignment
\[
(F,p)\longmapsto K_{\kappa_F}
\]
induces a bijection between reduced microlocal formal types of pointed boundary-linear vacua and polarized $A_\infty$ isomorphism classes of minimal pointed line algebras.
\end{corollary}

\begin{proof}
This is a restatement of Theorem~\ref{thm:line-mather-yau}.
\end{proof}

\subsection{What is stable and what is invisible}

The reduced microlocal formal type is obtained only after discarding linearly massive directions. The stabilization theorems of Section~\ref{sec:ts-knorrer} show that additional tangent-normal contractible blocks and pure-normal hyperbolic Clifford blocks are also invisible to the stable open sector.

\begin{corollary}[Stable invisibility of universal blocks; \ClaimStatusProvedHere]\label{cor:stable-invisibility}
The following operations do not change the stable open-sector content of a boundary model:
\begin{enumerate}[label=(\roman*),leftmargin=1.8em]
\item adding a tangent-normal hyperbolic pair $u v$ with one tangent and one normal leg;
\item adding a pure-normal hyperbolic pair $uv$ with both legs normal, after spinor stabilization.
\end{enumerate}
Hence the genuinely nontrivial local data of a pointed boundary-linear model are concentrated in the reduced Kuranishi map $\kappa$ and therefore in the polarized minimal pointed line algebra $K_\kappa$.
\end{corollary}

\begin{proof}
Part (i) is Theorem~\ref{thm:contractible-stab}. Part (ii) is Theorem~\ref{thm:clifford-stab}. The final assertion follows because Theorem~\ref{thm:boundary-kuranishi} removes the linearly massive directions explicitly and Theorem~\ref{thm:line-mather-yau} classifies what remains.
\end{proof}

\subsection{The fortified local triangle}

We may now summarize the exact local picture in one theorem.

\begin{theorem}[Fortified local bulk/boundary/line triangle; \ClaimStatusProvedHere]\label{thm:fortified-local-triangle}
For every pointed boundary-linear formal germ $(F,p)$ there is a canonical chain of identifications
\[
(F,p)
\rightsquigarrow
\cJ_{F,p}
\rightsquigarrow
K_{\mathrm{ex},F,p}
\simeq
K_{\kappa_F}
\rightsquigarrow
\HH^\bullet(K_{\kappa_F})
\simeq
\cO\bigl(\dCrit(W_{\eff,F})\bigr).
\]
Here:
\begin{itemize}[leftmargin=1.8em]
\item $\cJ_{F,p}$ is the one-step Jacobi coalgebra presenting the pointed local boundary equations;
\item $K_{\mathrm{ex},F,p}$ is the exact pointed line algebra;
\item $K_{\kappa_F}$ is the minimal pointed line algebra obtained by exact Kuranishi elimination;
\item and $\cO(\dCrit(W_{\eff,F}))$ is the completed local bulk algebra of the effective reduced potential.
\end{itemize}
The polarized $A_\infty$ isomorphism class of $K_{\kappa_F}$ determines the reduced microlocal formal type of $(F,p)$.
\end{theorem}

\begin{proof}
Combine Theorems~\ref{thm:exact-line}, \ref{thm:minimal-pointed-line}, \ref{thm:local-bulk-line}, and \ref{thm:line-mather-yau}.
\end{proof}

\begin{remark}[Swiss-cheese convolution algebra and 3d MC element]
\label{rem:forward-ref-sc-convolution}
The bulk-boundary-line triangle developed in this chapter is
encoded by a single algebraic object: the Swiss-cheese convolution
$L_\infty$-algebra $\gSC_T$
(Definition~\ref{def:sc-convolution}), whose canonical
Maurer--Cartan element $\alpha_T$
(Theorem~\ref{thm:3d-universal-mc}) recovers the bulk, boundary,
and line data as projections. These constructions are developed
in full in \S\ref{sec:sc-convolution-mc} of
Part~\ref{part:frontier}, following the same pattern as
Volume~I's modular convolution algebra $\gAmod$ and its universal
MC element~$\Theta_\cA$.
\end{remark}


\subsection{The modular homotopy type as complete invariant of the 3d theory}
\label{subsec:modular-homotopy-type}

The fortified local triangle of
Theorem~\ref{thm:fortified-local-triangle} exhausts what can be
proved exactly in the boundary-linear sector. But the bulk--boundary--line
triangle is only the genus-zero shadow of a richer structure: the
full modular homotopy type of the chiral algebra, which organises
every genus-$g$ obstruction into a single homotopical invariant.
The constructions that follow lift the local exact picture into
the modular convolution framework of Volume~I.

\begin{definition}[Modular homotopy type]
\label{def:modular-homotopy-type}
\index{modular homotopy type}
Let $\cA$ be a logarithmic $\SCchtop$-algebra.
The \emph{modular homotopy type} $\cM^{\mathrm{mod}}_\cA$ is the
simplicial Maurer--Cartan space
\begin{equation}\label{eq:modular-homotopy-type}
\cM^{\mathrm{mod}}_\cA
\;:=\;
\MC_\bullet\!\bigl(
 \Conv^{\mathrm{mod}}\!\bigl(\barB(\cA),\, \End(\cA)\bigr)
\bigr)
\end{equation}
equipped with the base-point $[\Theta_\cA]$ given by the
universal modular Maurer--Cartan element of Volume~I. Its
homotopy groups at this base-point encode:
\begin{itemize}[leftmargin=1.8em]
\item $\pi_0\bigl(\cM^{\mathrm{mod}}_\cA\bigr)$: isomorphism
 classes of chiral $\cA$-torsors (deformations of the identity
 bimodule);
\item $\pi_1\bigl(\cM^{\mathrm{mod}}_\cA, [\Theta_\cA]\bigr)$:
 automorphisms of the trivial torsor (the gauge group of $\cA$);
\item $\pi_n\bigl(\cM^{\mathrm{mod}}_\cA, [\Theta_\cA]\bigr)$
 for $n\ge 2$: higher symmetries and anomalies of the
 holomorphic--topological theory.
\end{itemize}
\end{definition}

\noindent
Here $\Conv^{\mathrm{mod}}\!\bigl(\barB(\cA),\End(\cA)\bigr)$
is the modular convolution $\Linf$-algebra obtained from the
stable-graph grading on $\gAmod$, and $\MC_\bullet$ denotes the
simplicial Maurer--Cartan functor in the sense of
Berglund~\cite{Berglund15} and Getzler~\cite{Getzler09}. The
base-point $\Theta_\cA$ is the canonical MC element constructed in
Volume~I (Theorem~A, the bar-cobar adjunction), whose projections recover
the chiral $\Ainf$ operations at every genus.

\begin{proposition}[Genus obstruction tower; \ClaimStatusProvedHere]
\label{prop:genus-tower-postnikov}
\index{modular homotopy type!obstruction tower}
The filtration of
$\Conv^{\mathrm{mod}}\!\bigl(\barB(\cA),\End(\cA)\bigr)$
by genus induces a tower
\begin{equation}\label{eq:postnikov-tower}
\cM^{\mathrm{mod}}_\cA
\;=\;
\varprojlim_{g}\, \cM^{\leq g}_\cA,
\end{equation}
where $\cM^{\leq g}_\cA$ is the simplicial MC space of the
genus-$\leq g$ truncation. The $g$-th layer, i.e.\ the
homotopy fiber of the restriction
$\cM^{\leq g}_\cA \to \cM^{\leq g-1}_\cA$,
has the following structure.
The homotopy fibre of $\cM^{\leq g}_\cA \to \cM^{\leq g-1}_\cA$ is a torsor for the Maurer--Cartan space of the genus-$g$ associated graded piece (an abelian $L_\infty$-algebra, hence a classifying space for the underlying chain complex). The \emph{obstruction} to the existence of a section (i.e., to extending a genus-$(g-1)$ MC element to genus~$g$) is the class of the right-hand side of~\eqref{eq:genus-g-obstruction} in $H^2$ of the associated graded; for $g = 1$ (where the quadratic sum is empty), this obstruction is exactly $\kappa(\cA) \cdot \omega_1$. For $g \geq 2$, the obstruction includes both the curvature term $\kappa(\cA) \cdot \omega_g$ and the quadratic terms $\sum [\Theta^{(g_1)}, \Theta^{(g_2)}]$ from lower genera.
\end{proposition}

\begin{proof}
The convolution algebra carries a complete exhaustive filtration
$F^{\leq g}$ indexed by the genus grading on stable graphs.
Completeness holds because $\mathrm{Conv}^{\mathrm{mod}}$ is defined as a product over genera: $\mathrm{Conv}^{\mathrm{mod}} = \prod_{g \geq 0} \mathrm{Conv}^{(g)}$, where each $\mathrm{Conv}^{(g)}$ is finite-dimensional (only finitely many stable graphs of genus $g$ with fixed degree contribute). The product topology is complete by definition, and $\mathrm{Conv}^{\mathrm{mod}} = \varprojlim_g \mathrm{Conv}^{\mathrm{mod}} / F^{\leq g}$.
The
associated graded at genus~$0$ is the $\Ainf$ bar complex; at
genus~$g\ge 1$, the graded piece consists of those structure maps
whose underlying graph has total genus~$g$. The MC equation
decomposes accordingly: a genus-$(g{-}1)$ MC element $\Theta^{\leq g-1}$
satisfies the truncated MC equation through genus~$g-1$. Its
extension to genus~$g$ requires solving
\begin{equation}\label{eq:genus-g-obstruction}
d_0\,\Theta^{(g)}
\;=\;
-\textstyle\sum_{\substack{g_1+g_2=g\\g_1,g_2<g}}
\tfrac{1}{2}[\Theta^{(g_1)},\Theta^{(g_2)}]
\;-\;\kappa(\cA)\cdot\omega_g,
\end{equation}
where $d_0$ is the genus-$0$ differential and the right-hand
side is determined by lower-genus data. The class of the
right-hand side in
$H^2\!\bigl(\Mbar_g,\,\End(\cA)\bigr)$ is
the obstruction to extending from genus~$g-1$ to genus~$g$.
For the universal element~$\Theta_\cA$, this obstruction is
$\kappa(\cA)\cdot\omega_1$ at genus~$1$ and, on the proved scalar
higher-genus lane, $\kappa(\cA)\cdot\omega_g$, by Volume~I Theorem~D
(leading coefficient).
At the level of simplicial sets, the filtered MC functor
converts a complete filtered $\Linf$-algebra into an inverse
limit of its truncated MC spaces
(Berglund~\cite[Prop.~3.5]{Berglund15}). The homotopy fiber of
each truncation map is identified with the MC space of the
associated graded piece, which is the one-step extension
controlled by the genus-$g$ obstruction.
\end{proof}

\begin{definition}[Quartic logarithmic contact class]
\label{def:quartic-log-contact-class}
\index{quartic logarithmic contact class|textbf}
\index{modular homotopy type!quartic contact class}
Let $\cA$ be a logarithmic $\SCchtop$-algebra and let
$\Theta^{\leq 3}_\cA$ denote the universal twisting element
determined by the binary and cubic structure of~$\cA$
(the genus-$0$, degree-$\leq 3$ component of~$\Theta_\cA$).
Define the \emph{quartic logarithmic contact class}
\begin{equation}\label{eq:quartic-log-contact}
\mathfrak{q}_4(\cA)
\;:=\;
\bigl[\pi_4\!\bigl(
 d\,\Theta^{\leq 3}_\cA
 + \tfrac{1}{2}[\Theta^{\leq 3}_\cA,\,\Theta^{\leq 3}_\cA]
\bigr)\bigr]
\;\in\;
H^2\!\bigl(\gr_4\, \Conv^{\mathrm{mod},\leq 4}(\barB(\cA),\End(\cA))\bigr),
\end{equation}
where $\pi_4$ projects to the weight-$4$ component in the
degree filtration. The class $\mathfrak{q}_4(\cA)$ is the
obstruction to extending the local genus-$0$ open/closed theory
from degree~$\leq 3$ across the first codimension-$2$ logarithmic
Fulton--MacPherson strata at degree~$4$.
\end{definition}

\begin{proposition}[Vanishing in the boundary-linear sector;
\ClaimStatusProvedHere]
\label{prop:quartic-contact-vanishes-BL}
\index{boundary-linear Landau--Ginzburg!quartic contact vanishing}
For every boundary-linear Landau--Ginzburg boundary algebra
$A_F = \bigl(k[\![x]\!] \otimes \Lambda(\eta),\; d\eta = F(x)\bigr)$,
\[
 \mathfrak{q}_4(A_F) = 0.
\]
\end{proposition}

\begin{proof}
In the boundary-linear sector the intrinsic boundary algebra is
strict: $m_n = 0$ for $n \geq 3$
(Corollary~\ref{cor:boundary-linear-strict}). Therefore the
twisting element $\Theta^{\leq 3}$ is determined entirely by the
strict differential $d$ and the binary product~$m_2$.

The weight-$4$ component of $d\,\Theta^{\leq 3}
+ \tfrac{1}{2}[\Theta^{\leq 3},\Theta^{\leq 3}]$ receives
contributions from codimension-$2$ boundary faces of the
genus-$0$ logarithmic cooperad $C^{\mathrm{oc,log}}_{\leq 4}$,
corresponding to two successive binary collisions or their
planted-forest refinements. Each such contribution is a
commutator $[m_2 \circ_i m_2,\ldots]$ involving the iterated
binary product.

Because $A_F$ is a strict commutative dg algebra:
\begin{enumerate}[label=(\alph*)]
\item $d^2 = 0$,
\item $d(ab) = (da)b + (-1)^{|a|} a\,(db)$ (Leibniz),
\item $(ab)c = a(bc)$ (associativity).
\end{enumerate}
The associativity identity~(c) forces all weight-$4$
planted-forest contributions to cancel in pairs: the two trees
contributing to the same codimension-$2$ corner produce
identical terms with opposite signs from the boundary orientation.
The Leibniz identity~(b) and $d^2 = 0$ eliminate the remaining
mixed differential--product terms. Therefore the total
obstruction class vanishes.
\end{proof}

\begin{remark}[Beyond the boundary-linear sector]
\label{rem:quartic-contact-beyond-BL}
For algebras with nontrivial higher operations $m_k$ ($k \geq 3$),
$\mathfrak{q}_4$ need not vanish. In the Virasoro family, the
quartic resonance class $\mathfrak{Q}^{\mathrm{contact}}_{\mathrm{Vir}}
= 10/[c(5c+22)]$ from the quartic-contact analysis of Volume~I is a
concrete instance
of a nonzero $\mathfrak{q}_4$. It is the first shadow that
distinguishes algebras with infinite tower depth (class~M) from
those with finite depth (classes~G, L, C): for Heisenberg
($\mathfrak{q}_4 = 0$, class~G), affine Kac--Moody
($\mathfrak{q}_4 = 0$, class~L), and $\beta\gamma$
($\mathfrak{q}_4 = 0$, class~C by stratum separation), the class
vanishes; for Virasoro and $\mathcal{W}_3$, it does not.
\end{remark}

In the computations below, the chiral algebra $\cA$ is the boundary algebra $A_b = \RHom_{\cC(J)}(b,b)$ of the vacuum boundary condition~$b$ (Theorem~\ref{thm:boundary-algebra-from-vacuum}).

\begin{computation}[Genus-$1$ obstruction for the Heisenberg algebra; \ClaimStatusProvedHere]
\label{comp:heisenberg-genus1-obstruction}%
\index{Heisenberg algebra!genus-1 obstruction}%
\index{modular homotopy type!Heisenberg genus-1 computation}%
We compute the genus-$1$ obstruction for the Heisenberg chiral
algebra $\cA = \cH_k$ step by step, grounding the tower of
Proposition~\ref{prop:genus-tower-postnikov} in an explicit
template calculation.

\smallskip
\noindent
\textbf{Step~1: genus-$0$ Maurer--Cartan element.}
The Heisenberg algebra has a single generator~$J$ with OPE
$J(z)\,J(w) \sim k/(z{-}w)^2$. All higher operations vanish:
$m_n = 0$ for $n \ge 3$
(Volume~I, Proposition~\ref{prop:rosetta-heisenberg-mk}).
The genus-$0$ MC element is therefore
\[
\Theta^{(0)} \;=\; m_2,
\qquad
m_2(J,J;\lambda) \;=\; k\lambda.
\]

\smallskip
\noindent
\textbf{Step~2: the genus-$1$ obstruction equation.}
Specialising~\eqref{eq:genus-g-obstruction} to $g = 1$:
\begin{equation}\label{eq:heisenberg-genus1-obstruction}
d_0\,\Theta^{(1)}
\;=\;
-\!\!\!\sum_{\substack{g_1 + g_2 = 1 \\ 0 < g_1 < 1}}
\tfrac{1}{2}[\Theta^{(g_1)},\Theta^{(g_2)}]
\;-\;
\kappa(\cH_k) \cdot \omega_1.
\end{equation}
The quadratic sum runs over pairs $(g_1,g_2)$ with
$g_1 + g_2 = 1$ and $0 < g_1 < 1$; there are no such pairs
of non-negative integers, so the sum is empty.

\smallskip
\noindent
\textbf{Step~3: evaluation of the curvature.}
Volume~I Theorem~D (leading coefficient) gives
$\kappa(\cH_k) = k$: the modular characteristic equals the
level. The obstruction equation reduces to
\begin{equation}\label{eq:heisenberg-genus1-reduced}
d_0\,\Theta^{(1)} \;=\; -k \cdot \omega_1,
\end{equation}
where
$\omega_1 = c_1(\lambda_{\mathrm{Hodge}})
\in H^2(\Mbar_{1,1},\,\mathbb{Z})$
is the first Chern class of the Hodge bundle
$\lambda_{\mathrm{Hodge}}$ on $\Mbar_{1,1}$.

\smallskip
\noindent
\textbf{Step~4: cohomology of the genus-$1$ graph complex.}
The obstruction to solving~\eqref{eq:heisenberg-genus1-reduced}
lies in
$H^1\!\bigl(C^{(1)}_{\cH_k},\, d_0\bigr)$,
the genus-$1$ graph complex with binary vertices
(since $m_{n \ge 3} = 0$). A graph of genus~$1$ with only
binary vertices has a single topological type: one vertex with
two half-edges joined into a single self-loop edge. Write
$\Gamma_{\mathrm{irr}}$ for this graph (the irreducible boundary
graph on $\Mbar_{1,1}$, corresponding to the boundary divisor
$\Delta_{\mathrm{irr}} \subset \Mbar_{1,1}$). The complex
$(C^{(1)},d_0)$ is therefore one-dimensional, spanned by
$\Gamma_{\mathrm{irr}}$.

\smallskip
\noindent
\textbf{Step~5: explicit solution.}
The solution is
\begin{equation}\label{eq:heisenberg-theta1}
\Theta^{(1)}
\;=\;
-k \cdot \Gamma_{\mathrm{irr}} \otimes \Lambda,
\end{equation}
where
$\Lambda \in H^1(\Mbar_{1,1})$ is the fundamental class
of the boundary divisor $\Delta_{\mathrm{irr}}$. To verify:
the self-sewing operation on $\Gamma_{\mathrm{irr}}$ contracts
the propagator through the self-loop, yielding
\[
d_0\!\bigl(\Gamma_{\mathrm{irr}} \otimes \Lambda\bigr)
\;=\;
\omega_1
\]
by the clutching formula for $\Mbar_{1,1}$: the boundary
$\Delta_{\mathrm{irr}}$ contributes
$c_1(\lambda_{\mathrm{Hodge}})$ through the self-sewing map
$\Mbar_{0,2+2} \to \Delta_{\mathrm{irr}} \hookrightarrow
\Mbar_{1,1}$. Therefore
$d_0\,\Theta^{(1)} = -k \cdot \omega_1$,
matching~\eqref{eq:heisenberg-genus1-reduced}.

\smallskip
\noindent
\textbf{Step~6: free energy.}
The genus-$1$ free energy is the scalar obtained by evaluating
$\Theta^{(1)}$ on the moduli space. Since
$\Theta^{(1)} = -k \cdot \Gamma_{\mathrm{irr}} \otimes \Lambda$
and $\int_{\Mbar_{1,1}} \omega_1 = 1/24$ (Mumford's formula):
\begin{equation}\label{eq:heisenberg-F1}
F_1(\cH_k)
\;=\;
-k\,\log\,\eta(\tau),
\end{equation}
where $\eta(\tau) = q^{1/24}\prod_{n=1}^\infty (1 - q^n)$
is the Dedekind eta function and $q = e^{2\pi i \tau}$.
This recovers the standard one-loop free energy of the
free boson at level~$k$ at fixed modular parameter~$\tau$.
(The factor $1/24$ from Mumford's integral
$\int_{\Mbar_{1,1}} \lambda_1 = 1/24$ enters the
\emph{integrated} partition function
$Z_1 = \int e^{F_1}\,d\mu$, not the local free
energy $F_1(\tau)$ at fixed~$\tau$.)

\smallskip
\noindent
\textbf{Step~7: shadow depth verification.}
The solution $\Theta^{(1)}$ is determined entirely by the
curvature $\kappa(\cH_k) = k$ and the binary operation $m_2$:
no independent genus-$1$ data is required. Consequently,
$d_{\mathrm{alg}} = 0$, confirming the class~$\mathbf{G}$
(Gaussian) designation of
Proposition~\ref{prop:shadow-depth}. At every genus
$g \ge 1$, the MC element $\Theta^{(g)}$ is uniquely determined
by the self-sewing recursion applied to $\kappa = k$, with no
free parameters.
\end{computation}

\begin{proposition}[Shadow depth as structural complexity; \ClaimStatusProvedHere]
\label{prop:shadow-depth}
\index{shadow depth}
The \emph{shadow depth} $d(\cA)$ is the minimal integer $r$
such that the $r$-th truncation $\Theta^{\leq r}_\cA$ of the
universal MC element determines $\Theta_\cA$ up to
quasi-isomorphism of $\SCchtop$-algebras. The four depth
classes are:
\begin{description}[leftmargin=2em]
\item[\textbf{G} (Gaussian, $d_{\mathrm{alg}}=0$):]
 Free algebras. The bar complex is cofree, all higher
 operations vanish, and the MC element is the linear term alone.
 The $R$-matrix is trivial (or the exponential of a scalar).
 \emph{Example:} Heisenberg $\cH_k$.
\item[\textbf{L} (Lie, $d_{\mathrm{alg}}=1$):]
 Quadratic algebras. The bar differential has only binary
 operations; the quantum $R$-matrix
 $R(z) = \mathrm{Pexp}\!\bigl(\int r(w)\,dw\bigr)$
 satisfies the full quantum Yang--Baxter equation, but is
 \emph{computable from $m_2$ alone}: it is determined
 entirely by the classical $r$-matrix (binary OPE data).
 \emph{Example:} affine $\hat{\mathfrak{g}}_k$ at generic level.
\item[\textbf{C} (Contact, $d_{\mathrm{alg}}=2$):]
 Cubic algebras. The quantum $R$-matrix satisfies the full
 quantum Yang--Baxter equation (as it must for any
 $\SCchtop$-algebra), but requires independent ternary data
 ($m_3$) beyond the classical $r$-matrix for its determination.
 \emph{Example:} $\beta\gamma$ system.
\item[\textbf{M} (Mixed, $d_{\mathrm{alg}}=\infty$):]
 No finite truncation of the bar data determines the
 $R$-matrix. The full perturbative series of
 higher operations contributes at every degree.
 \emph{Example:} Virasoro at generic~$c$; $\cW_3$ at non-generic central charge.
\end{description}
\end{proposition}

\begin{proof}
The shadow depth $d(\cA)$ is defined by the degree filtration on
the bar coderivation: it is the smallest $r$ such that all MC
components $\Theta^{(n)}_\cA$ for $n>r$ are determined
(via the Stasheff system) by those for $n\le r$. For a free
algebra, the bar coalgebra carries no differential beyond the
linear counit, so $d=0$. For a quadratic algebra, the bar
differential is determined by $m_2$, so $d=1$; the quantum
$R$-matrix is then computable as $R(z)=\mathrm{Pexp}(\int r(w)\,dw)$
from the classical $r$-matrix alone, and it satisfies the full
quantum Yang--Baxter equation (the degree-$2$ MC equation in the
$E_1$-modular convolution algebra already encodes the quantum
relation).
For a cubic algebra, the degree-$3$ Stasheff relation (expressing $m_3$ as a homotopy for the failure of $m_2$ to be
associative, mediated by $m_1$) is not automatically satisfied
by $m_2$ alone, so $m_3$ is
independent data and $d=2$. If the algebra has essential
operations at every degree (as happens for $\cW_3$ at
non-generic central charge, where the Kuranishi map $\kappa$ has
Taylor coefficients at all orders), then no finite truncation
suffices and $d=\infty$.
\end{proof}

\begin{computation}[Shadow depth of $\hat{\mathfrak{sl}}_2$ at generic level; \ClaimStatusProvedHere]
\label{comp:shadow-depth-sl2}
\index{shadow depth!affine $\hat{\mathfrak{sl}}_2$}

We derive that $\hat{\mathfrak{sl}}_2$ at generic level $k$
(i.e.\ $k\neq -2$) has shadow depth $d_{\mathrm{alg}}=1$
(class~L) by showing: (i)~$m_2$ is independent data,
(ii)~all $m_{n\ge 3}$ are determined by $m_2$ via the Stasheff
relations.

\medskip
\noindent\textbf{Step 1: OPE data.}
The affine vertex algebra $V_k(\mathfrak{sl}_2)$ has
generators $J^a(z)$, $a\in\{+,-,3\}$, with OPE
\begin{equation}\label{eq:sl2-ope}
J^a(z)\, J^b(w) \;\sim\;
\frac{k\,\delta^{ab}}{(z-w)^2}
\;+\;
\frac{f^{ab}{}_{c}\, J^c(w)}{z-w},
\end{equation}
where $f^{ab}{}_{c}$ are the $\mathfrak{sl}_2$ structure
constants (normalised so that $[e,f]=h$, $[h,e]=2e$,
$[h,f]=-2f$) and $\delta^{ab}$ is the Killing form
normalised to $\delta^{ab}=\tfrac{1}{2}\operatorname{tr}
(\operatorname{ad}_{J^a}\circ\operatorname{ad}_{J^b})$.

\medskip
\noindent\textbf{Step 2: Bar differential at degree 2.}
The bar differential at degree~$2$ encodes both the
double-pole (level) and simple-pole (structure constant)
data in a single $\lambda$-bracket:
\begin{equation}\label{eq:sl2-m2}
m_2(J^a, J^b;\lambda)
\;=\;
k\,\delta^{ab}\,\lambda
\;+\;
f^{ab}{}_{c}\, J^c.
\end{equation}
The first term is the Killing-form contribution (the
central extension); the second is the Lie bracket.
This is non-trivial data: neither the level $k$ nor the
structure constants are determined by lower-degree
operations, so $d_{\mathrm{alg}}\ge 1$.

\medskip
\noindent\textbf{Step 3: Vanishing of independent
 higher operations.}
At degree~$3$, the bar differential computes the triple
OPE. For $\hat{\mathfrak{sl}}_2$ at generic level
$k\neq -h^\vee = -2$, the Sugawara construction is
non-degenerate: the Sugawara energy-momentum tensor
$T^{\mathrm{Sug}}(z)=\frac{1}{2(k+2)}\sum_a
{:}\!J^a J_a\!{:}(z)$ exists and gives a conformal vector.
The one-loop exactness theorem
(Theorem~\ref{thm:one-loop-koszul}) then shows that the
BV-BRST differential of the associated 3d HT theory has
no quantum corrections beyond one loop. Concretely,
all Feynman diagrams contributing to $m_3$ are trees with
binary vertices, already accounted for by
iterating $m_2$. The degree-$3$ Stasheff relation
(Definition~\ref{def:ainfty-relations} at $n=3$), which expresses the homotopy associativity of $m_2$ controlled by $m_3$ with the differential $m_1$ acting by Leibniz, determines $m_3$ uniquely as a chain homotopy for the
failure of $m_2$ to be strictly associative. Since
$m_2\circ m_2$ is already known from \eqref{eq:sl2-m2},
$m_3$ carries no independent information. By induction on
degree: if $m_2,\ldots,m_{n-1}$ are all determined by
$m_2$, the degree-$n$ Stasheff relation expresses $m_n$ in
terms of compositions of $m_2,\ldots,m_{n-1}$, hence in
terms of $m_2$ alone. Therefore the bar coderivation is
determined by its degree-$2$ component, giving
$d_{\mathrm{alg}}=1$.

\medskip
\noindent\textbf{Step 4: Classical $r$-matrix.}
The classical $r$-matrix extracted from $m_2$ is
\begin{equation}\label{eq:sl2-r-matrix}
r(z) \;=\; \frac{k\,\Omega}{z}
\;=\; \frac{k}{z}\sum_{a} J^a \otimes J_a,
\end{equation}
where $\Omega = \sum_a J^a\otimes J_a$ is the quadratic
Casimir element of $\mathfrak{sl}_2$ and $k$ is the level
of $\widehat{\mathfrak{sl}}_2$. One checks directly
that $r(z)$ satisfies the classical Yang--Baxter equation
with spectral parameter:
\[
[r_{12}(z_1{-}z_2),\, r_{13}(z_1{-}z_3)]
+ [r_{12}(z_1{-}z_2),\, r_{23}(z_2{-}z_3)]
+ [r_{13}(z_1{-}z_3),\, r_{23}(z_2{-}z_3)]
= 0,
\]
which is the Jacobi identity for $\mathfrak{sl}_2$ divided
by $(z_1{-}z_2)(z_1{-}z_3)(z_2{-}z_3)$. The quantum
$R$-matrix
\begin{equation}\label{eq:sl2-R-matrix}
R(z) \;=\; \mathcal{P}\exp\!\Bigl(\int_0^z
r(w)\,dw\Bigr)
\end{equation}
satisfies the full quantum Yang--Baxter equation
(Theorem~\ref{thm:YBE}), and is determined entirely by
$r(z)$, confirming class~L.

\medskip
\noindent\textbf{Step 5: Genus-$1$ obstruction.}
%: kappa(V_k(sl_2)) = dim(sl_2)(k+h^v)/(2h^v) = 3(k+2)/4; NOT k
The modular characteristic is $\kappa(\hat{\mathfrak{sl}}_2)
= \frac{3(k+2)}{4}
= \frac{\dim(\mathfrak{sl}_2)(k+h^\vee)}{2h^\vee}$, and the genus-$1$ obstruction takes the form
\begin{equation}\label{eq:sl2-genus1}
\Theta^{(1)} \;=\;
-\frac{3(k+2)}{4} \cdot \Gamma_{\text{self-loop}} \otimes \Lambda,
\end{equation}
where $\Gamma_{\text{self-loop}}$ is the unique genus-$1$
graph with a single vertex and one loop, and $\Lambda$
carries an $\mathfrak{sl}_2$ Casimir insertion at the
vertex. (The sign is fixed by the obstruction equation
$d_0\,\Theta^{(1)} = -\frac{3(k+2)}{4}\cdot\omega_1$ together with
$d_0(\Gamma_{\text{self-loop}} \otimes \Lambda) = \omega_1$.)
The genus-$1$ free-energy coefficient for
$\hat{\mathfrak{sl}}_2$ at level~$k$ is
\[
F_1(\hat{\mathfrak{sl}}_2,\, k)
\;=\;
\frac{\kappa(\hat{\mathfrak{sl}}_2)}{24}
\;=\;
\frac{3(k+2)}{4\cdot 24}
\;=\;
\frac{k+2}{32}.
\]
Here $\dim(\mathfrak{sl}_2)=3$ and $h^\vee=2$, so the affine
Kac--Moody formula for $\kappa$ gives the stated value. At
$k=0$ one gets $F_1 = 1/16$, while at the critical level
$k=-2$ the coefficient vanishes with $\kappa$. The structural form matches the
Heisenberg genus-$1$ obstruction
(Computation~\ref{comp:heisenberg-genus1-obstruction}),
with the scalar self-loop now carrying the Casimir
$\Omega = \sum_a J^a \otimes J_a$ in place of the identity.
The class~L property ensures that
$\Theta^{(1)}$, like all higher-genus
obstructions, is computable from $m_2$ alone.
\end{computation}

\begin{theorem}[Complementarity as holographic dictionary; \ClaimStatusProvedHere]
\label{thm:complementarity-holographic-dictionary}
\index{complementarity!holographic dictionary}
\index{holographic dictionary!via complementarity}
Let $\cA$ be a chirally Koszul logarithmic $\SCchtop$-algebra.
At each genus $g\ge 1$, the modular obstruction space
decomposes as
\begin{equation}\label{eq:complementarity-decomposition}
Q_g(\cA) \;\oplus\; Q_g(\cA^!)
\;=\;
H^\bullet\!\bigl(\Mbar_g,\, \cZ_\cA\bigr),
\end{equation}
where:
\begin{itemize}[leftmargin=1.8em]
\item $\cZ_\cA$ is the \emph{center constructible sheaf} on $\Mbar_g$: its fiber at a smooth curve $C$ is the genus-$g$ modular deformation complex $\mathrm{Def}^{(g)}(\cA|_C)$ of the chiral algebra $\cA$ restricted to $C$ (when $\cA$ is smooth, $\cZ_\cA$ is a local system);
\item $Q_g(\cA)$ is the \emph{deformation sector}
 (boundary-accessible obstructions);
\item $Q_g(\cA^!)$ is the \emph{obstruction sector}
 (boundary-inaccessible obstructions).
\end{itemize}
The Verdier involution
$\sigma\colon H^\bullet(\Mbar_g,\cZ_\cA)\to H^\bullet(\Mbar_g,\cZ_\cA)$,
with $\sigma^2=\mathrm{id}$, exchanges
the two sectors:
$\sigma\bigl(Q_g(\cA)\bigr)=Q_g(\cA^!)$.
This is the holographic dictionary: the bulk genus-$g$ Hilbert
space $H^\bullet(\Mbar_g,\cZ_\cA)$ splits into
boundary-accessible and boundary-inaccessible sectors, and
Koszul duality is the operation that transposes the split.
\end{theorem}

\begin{construction}[The residue pairing on the modular obstruction space; \ClaimStatusProvedHere]
\label{constr:residue-pairing}
\index{residue pairing!modular obstruction space}
\index{shifted symplectic pairing!construction}

We construct the $(-1)$-shifted symplectic pairing on the genus-$g$
modular obstruction space
$H^\bullet(\Mbar_g,\,\cZ_\cA)$
invoked in Theorem~\ref{thm:complementarity-holographic-dictionary}.

\medskip
\noindent
\textit{Step~1: The bilinear form.}\enspace
Let $\cA$ be a chirally Koszul logarithmic $\SCchtop$-algebra.
The genus-$g$ deformation complex
$\mathfrak{L}_g = \barB^{(g)}(\cA)[1]$
carries a bilinear form
\begin{equation}\label{eq:residue-pairing-definition}
\langle \alpha,\, \beta \rangle_{-1}
\;:=\;
\int_{\Mbar_{g,n}}
\mathrm{Res}\!\bigl(\alpha \cdot \beta\bigr),
\end{equation}
defined as follows. For bar chains
$\alpha = [s^{-1}a_1 \,|\, \cdots \,|\, s^{-1}a_p]$
and
$\beta = [s^{-1}b_1 \,|\, \cdots \,|\, s^{-1}b_q]$,
the product $\alpha \cdot \beta$ is defined via the coalgebra
diagonal $\Delta\colon \barB(\cA) \to \barB(\cA) \otimes \barB(\cA)$
and the Serre duality pairing on $\cA$:
\begin{equation}\label{eq:residue-pairing-matchings}
\mathrm{Res}(\alpha \cdot \beta)
\;=\;
\sum_{\sigma}
\prod_{i=1}^{p}
\langle a_i,\, b_{\sigma(i)} \rangle_{\mathrm{Serre}}
\;\cdot\;
\omega_{p+q},
\end{equation}
where the sum runs over all matchings $\sigma$ of the bar
factors and $\omega_{p+q}$ is the weight form on
$\FM_{p+q}(\C)$. The matching sum arises from the components of
the iterated coproduct
$\Delta^{(p+q-1)}\colon \barB^{(g)}(\cA)
\to \barB(\cA)^{\otimes(p+q)}$
composed with the Serre trace
$\langle\,\cdot\,,\,\cdot\,\rangle_{\mathrm{Serre}}
\colon \cA \otimes \cA \to \kk$
on matched pairs.

\medskip
\noindent
\textit{Step~2: Degree verification.}\enspace
The pairing has total degree $-1$. The bar desuspension
contributes $-p$ from the $p$ factors $s^{-1}$ in $\alpha$ and
$-q$ from the $q$ factors in $\beta$, for a total of $-(p+q)$.
The Serre pairing on a curve ($\dim_\C X = 1$) contributes
$+1$ per matched pair, giving $+\min(p,q)$; when $p = q$ this is
$+p$. The weight form $\omega_{p+q}$ is the top-degree
logarithmic form on $\FM_{p+q}(\C)$, contributing
$+(p+q-1)$. For the pairing to be nonzero we need $p = q$
(otherwise the matching sum vanishes), giving total degree:
\begin{equation}\label{eq:residue-pairing-degree}
-(p+q) + p + (p+q-1) \;=\; p - 1.
\end{equation}
But the integral over $\Mbar_{g,n}$ consumes the
$\Mbar_g$-degree $p - 1$ (the relevant component of the
modular convolution complex has $\Mbar_g$-weight $p-1$), leaving
a residual degree of
\[
(p-1) - (p-1) + (-1) \;=\; -1,
\]
where the final $-1$ comes from the cohomological shift
$\mathfrak{L}_g = \barB^{(g)}(\cA)[1]$. Thus the pairing
$\langle\,\cdot\,,\,\cdot\,\rangle_{-1}$ has degree $-1$ as
required for a $(-1)$-shifted symplectic form.

\medskip
\noindent
\textit{Step~3: Nondegeneracy.}\enspace
On the Koszul locus (chirally Koszul $\cA$), the bar complex
is a resolution: $\barB(\cA) \xrightarrow{\sim} \cA^{!,\vee}$
by Volume~I Theorem~B. Serre duality on the curve $X$ provides
a nondegenerate pairing $\cA \otimes \cA^{!} \to \kk$, and the
Koszul resolution property promotes this to a chain-level
quasi-isomorphism
\begin{equation}\label{eq:residue-pairing-nondeg}
\barB^{(g)}(\cA)
\;\xrightarrow{\;\sim\;}
\barB^{(g)}(\cA)^\vee[-1].
\end{equation}
Concretely: the map sends a bar chain $\alpha$ to the functional
$\beta \mapsto \langle \alpha,\beta\rangle_{-1}$. This is a
quasi-isomorphism because (i) the Serre pairing is
nondegenerate on $\cA$ (by assumption: $\cA$ is a
factorization algebra on a smooth curve), (ii) the bar complex
resolves $\cA^{!,\vee}$ (Koszul property), and (iii) the weight
form $\omega_{p+q}$ pairs the top and bottom strata of
$\FM_{p+q}(\C)$ nondegenerately (Poincar\'{e} duality for manifolds
with corners, applied to the FM compactification). On
cohomology, this gives a nondegenerate pairing on
$H^\bullet(\Mbar_g,\,\cZ_\cA)$.

\medskip
\noindent
\textit{Step~4: Closedness.}\enspace
The pairing is closed (i.e., compatible with the differential):
$\langle d\alpha,\,\beta\rangle_{-1}
+ (-1)^{|\alpha|}\langle\alpha,\,d\beta\rangle_{-1} = 0$.
This follows from the fact that the bar differential $d_B$ is a
\emph{coderivation} for the coalgebra structure $\Delta$, and the
Serre pairing satisfies the invariance condition
$\langle d_\cA(a),\,b\rangle_{\mathrm{Serre}}
= -(-1)^{|a|}\langle a,\,d_\cA(b)\rangle_{\mathrm{Serre}}$.
The Stokes formula on $\FM_{p+q}(\C)$ supplies the boundary
contributions, which cancel by the Arnold relations (exactness of
the weight forms on boundary strata).

\medskip
\noindent
\textit{Step~5: Isotropy of the Lagrangians.}\enspace
The bar Lagrangian
$Q_g(\cA) = \im\!\bigl(\barB^{(g)}(\cA)
\to H^\bullet(\Mbar_g,\,\cZ_\cA)\bigr)$
is isotropic: for $\alpha,\,\beta \in Q_g(\cA)$, we have
$\langle\alpha,\,\beta\rangle_{-1} = 0$. To see this, observe
that both $\alpha$ and $\beta$ lift to bar chains. In the
matching sum~\eqref{eq:residue-pairing-matchings}, the Serre
pairing $\langle a_i,\, b_{\sigma(i)}\rangle_{\mathrm{Serre}}$
pairs two elements that both come from $\cA$ (the bar side).
Under the identification $\barB(\cA) \simeq \cA^{!,\vee}$, the
Serre pairing between two elements of $\cA^{!,\vee}$ factors
through
$\cA^{!,\vee} \otimes \cA^{!,\vee} \to (\cA^! \otimes \cA^!)^\vee$,
but the image lands in the \emph{symmetric} part of
$(\cA^! \otimes \cA^!)^\vee$ (the bar complex is cofree on the
cosuspension, and the coproduct is cocommutative on cohomology).
Since the residue pairing has odd degree $(-1)$, it is
\emph{antisymmetric} on the symmetric part, hence vanishes. An
analogous argument shows that the cobar Lagrangian
$Q_g(\cA^!) = \im\!\bigl(\barB^{(g)}(\cA^!)
\to H^\bullet(\Mbar_g,\,\cZ_\cA)\bigr)$
is also isotropic. Together with the dimension count
$\dim Q_g(\cA) + \dim Q_g(\cA^!)
= \dim H^\bullet(\Mbar_g,\,\cZ_\cA)$
from Volume~I Theorem~C, this establishes that $Q_g(\cA)$ and
$Q_g(\cA^!)$ are complementary Lagrangians.
\end{construction}

\begin{construction}[Cyclic structure on the open-sector category; \ClaimStatusProvedHere]
\label{constr:cyclic-structure}
The $(-1)$-shifted symplectic pairing of Construction~\ref{constr:residue-pairing} induces a \emph{cyclic $A_\infty$ structure} on the boundary algebra $A_b = \End_\cC(b)$.

\textbf{Step 1: The cyclic pairing.} Define the nondegenerate bilinear form
\[
\langle a, b \rangle_{\mathrm{cyc}} \;:=\; \Res_{z=0}\, \langle a(z),\, b(0) \rangle_{\mathrm{Serre}} \cdot dz
\]
on $A_b$, where $\langle -,- \rangle_{\mathrm{Serre}}$ is the Serre duality pairing on the curve $X$. This pairing has degree $-1$ (from the $(-1)$-shifted symplectic structure on the self-intersection $\cL \times_{\cM^h} \cL$).

\textbf{Step 2: Cyclic $A_\infty$ identity.} The pairing is cyclic with respect to all $A_\infty$ operations:
\begin{equation}\label{eq:cyclic-Ainfty}
\langle m_k(a_1, \ldots, a_k;\, \lambda_1, \ldots, \lambda_{k-1}),\, a_0 \rangle_{\mathrm{cyc}} \;=\; (-1)^{\dagger}\, \langle m_k(a_0, a_1, \ldots, a_{k-1};\, \lambda_0, \ldots, \lambda_{k-2}),\, a_k \rangle_{\mathrm{cyc}}
\end{equation}
where $(-1)^\dagger = (-1)^{|a_0|(|a_1|+\cdots+|a_k|)}$ is the Koszul sign for cyclic permutation of graded inputs, and the spectral parameters rotate correspondingly: $\lambda_0 = -(\lambda_1 + \cdots + \lambda_{k-1} + \partial)$ (the momentum-conservation constraint). Here $|a_i|$ denotes the cohomological degree in the \emph{unsuspended} (original) grading. In the bar complex, the desuspended elements $s^{-1}a_i$ have bar degree $\overline{a}_i = |a_i| - 1$. The sign $(-1)^\dagger = (-1)^{|a_0|(|a_1|+\cdots+|a_k|)}$ is equivalent to the standard Kontsevich--Soibelman sign $(-1)^{\overline{a}_0(\overline{a}_1+\cdots+\overline{a}_k)}$ via the identity $|a_0|(|a_1|+\cdots+|a_k|) \equiv \overline{a}_0(\overline{a}_1+\cdots+\overline{a}_k) + \overline{a}_0 \cdot k + (|a_1|+\cdots+|a_k|) \pmod{2}$, with the additional terms absorbed into the definition of the pairing (the degree $(-1)$ of $\langle -,- \rangle_{\mathrm{cyc}}$ contributes $(-1)^{\overline{a}_0}$ from commuting $a_0$ past the odd pairing, and the $A_\infty$ degree $|m_k| = 2-k$ contributes $(-1)^{k-1}$, cancelling the extra terms).

\textbf{Step 3: Proof from Serre duality.} The cyclic identity follows from the $\SCchtop$-equivariance of Serre duality: the pairing $\langle -,- \rangle_{\mathrm{Serre}}$ is compatible with the factorization structure because Serre duality on $X$ intertwines the chiral product $\mu$ with its transpose. The cyclic rotation of inputs corresponds to the rotation of marked points on the boundary circle $S^1_p$ of the bordified curve, which is a symmetry of the Serre pairing. The spectral parameter rotation is the sesquilinearity constraint (Lemma~\ref{lem:sesquilinearity_explicit}) applied to the rotated configuration.

\textbf{Step 4: The trace map.} The cyclic structure determines a trace:
\[
\Tr_\cC \colon \HH_\ast(\cC) \;\longrightarrow\; \kk,
\]
defined on the Hochschild homology of $\cC$ by contracting with the cyclic pairing. For a Hochschild cycle $\sum a_0 \otimes a_1 \otimes \cdots \otimes a_n$ (a cyclic bar chain), the trace is:
\[
\Tr_\cC(a_0 \otimes \cdots \otimes a_n) \;=\; \langle m_n(a_1, \ldots, a_n),\, a_0 \rangle_{\mathrm{cyc}}.
\]
The cyclic identity~\eqref{eq:cyclic-Ainfty} ensures this is well-defined on cyclic equivalence classes.
\end{construction}

\begin{remark}[Cyclic structure and modularity]
\label{rem:cyclic-modularity}
The cyclic structure is the datum that makes modular gluing possible: genus-$g$ amplitudes arise from composing $g$ clutching operations with the trace $\Tr_\cC$. The modular characteristic $\kappa(\cA)$ measures the obstruction to the trace being a \emph{chain map} (i.e., compatible with the bar differential at genus $g \geq 1$). When $\kappa = 0$, the trace is a chain map and the theory is ``anomaly-free''; when $\kappa \neq 0$, the trace acquires a curvature correction $\kappa \cdot \omega_g$, which is the conformal anomaly.
\end{remark}

\begin{proof}
Volume~I Theorem~C (complementarity) establishes that the bar
and cobar complexes provide complementary Lagrangian
decompositions of the modular obstruction space at each genus.
The shifted symplectic structure invoked here is the BV
structure: a $(-1)$-shifted symplectic form on the formal moduli
problem $\MC(\mathfrak{L}_g)$ (not on the cohomology group
$H^\bullet(\Mbar_g,\cZ_\cA)$ itself). The Lagrangian
decomposition is the splitting
$\mathfrak{L}_g = \mathfrak{L}_g^+ \oplus \mathfrak{L}_g^-$ into bar and cobar eigenspaces
of the Verdier involution, with isotropy and nondegeneracy
verified at the chain level.

Concretely, the residue pairing on the modular convolution
complex descends from this $(-1)$-shifted symplectic form, and
the bar image
$\im\!\bigl(\barB(\cA)\to H^\bullet(\Mbar_g,\cZ_\cA)\bigr)$
forms an isotropic subspace whose rank equals half the total
dimension by the complementarity theorem. The cobar image
$\im\!\bigl(\Omegach(\barB(\cA))\to H^\bullet(\Mbar_g,\cZ_\cA)\bigr)$
provides the complementary Lagrangian.

For chirally Koszul $\cA$, the Koszul involution $\cA\leftrightarrow\cA^!$
exchanges bar and cobar:
$\barB(\cA)\simeq\cA^{!,\vee}$, so the
cobar Lagrangian for $\cA$ is the bar Lagrangian for $\cA^!$.
Setting $Q_g(\cA):=\im(\barB(\cA))$ and
$Q_g(\cA^!):=\im(\barB(\cA^!))$, the complementarity theorem
gives the direct-sum decomposition~\eqref{eq:complementarity-decomposition}.

The Verdier involution $\sigma$ on the center constructible sheaf
$\cZ_\cA$ induces an involution
$\sigma\colon H^\bullet(\Mbar_g,\cZ_\cA)\to H^\bullet(\Mbar_g,\cZ_\cA)$
with $\sigma^2=\mathrm{id}$, acting as the $(-1)$-shifted symplectic
duality. Its exchange of the two
complementary Lagrangians is precisely the statement
$\sigma(Q_g(\cA))=Q_g(\cA^!)$.

The holographic interpretation follows: elements of $Q_g(\cA)$
are visible to the boundary algebra $\Bbound$ because they
lie in the image of the bar construction of $\cA$, which
computes $\Zder(\Bbound)$ via the bulk--boundary
theorem. Elements of $Q_g(\cA^!)$ require passage through
the Koszul dual and are therefore inaccessible from the
boundary alone.
\end{proof}

\begin{computation}[Complementarity decomposition for $\cH_k$ at genus~$1$;
\ClaimStatusProvedHere]
\label{comp:complementarity-heisenberg-g1}
\index{complementarity!Heisenberg!genus 1}
\index{Heisenberg algebra!complementarity at genus 1}

We exhibit the Lagrangian decomposition
$Q_1(\cH_k) \oplus Q_1(\cH_k^!) = H^\bullet(\Mbar_{1,1},\, \cZ_{\cH_k})$
of Theorem~\ref{thm:complementarity-holographic-dictionary}
explicitly for the Heisenberg algebra at genus~$1$.
We work on $\Mbar_{1,1}$ (one marked point), where the Hodge
bundle is nontrivial and the complementarity structure is
cleanly visible.

\medskip
\noindent
\textit{Step~1: Cohomology of $\Mbar_{1,1}$.}\enspace
As a Deligne--Mumford stack, $\Mbar_{1,1}$ has the homotopy
type of a weighted projective line. Its rational cohomology
ring is
\begin{equation}\label{eq:complementarity-Mbar11-cohomology}
H^\bullet(\Mbar_{1,1},\,\kk)
\;=\;
\kk[\omega_1]/(\omega_1^2),
\qquad
|\omega_1| = 2,
\end{equation}
where $\omega_1 = c_1(\mathbb{E})$ is the first Chern class of
the Hodge bundle $\mathbb{E} \to \Mbar_{1,1}$. In particular,
$H^0 = \kk \cdot 1$, $H^1 = 0$, and
$H^2 = \kk \cdot \omega_1$, so $\dim H^\bullet = 2$.

\medskip
\noindent
\textit{Step~2: The center constructible sheaf.}\enspace
The Heisenberg algebra $\cH_k$ has a single generator~$J$ of
conformal weight~$1$ and OPE
$J(z)\,J(w) \sim k/(z-w)^2$.
The modular deformation complex at genus~$1$ is
$\mathrm{Def}^{(1)}(\cH_k) = \kk$, one-dimensional, spanned by
the modular characteristic $\kappa(\cH_k) = k$
(Volume~I Theorem~D; see also
equation~\eqref{eq:heisenberg-genus1-reduced}). Since the
fiber is constant over $\Mbar_{1,1}$, the center constructible
sheaf is the constant sheaf:
\[
\cZ_{\cH_k}
\;\cong\;
\underline{\kk}_{\Mbar_{1,1}}.
\]
Consequently, the total modular obstruction space is
\begin{equation}\label{eq:complementarity-total-space}
H^\bullet(\Mbar_{1,1},\, \cZ_{\cH_k})
\;=\;
H^\bullet(\Mbar_{1,1},\,\kk)
\;=\;
\kk \cdot 1 \;\oplus\; \kk \cdot \omega_1,
\end{equation}
a two-dimensional $\kk$-vector space concentrated in
degrees~$0$ and~$2$.

\medskip
\noindent
\textit{Step~3: The $(-1)$-shifted symplectic pairing.}\enspace
The Verdier pairing on $H^\bullet(\Mbar_{1,1},\,\cZ_{\cH_k})$
is the $(-1)$-shifted intersection form
\[
\langle \alpha,\, \beta \rangle_{-1}
\;=\;
\int_{\Mbar_{1,1}} \alpha \cup \beta
\;\in\; \kk,
\]
shifted so that elements of complementary degree pair
nondegenerately. On the two-dimensional space
$\kk \cdot 1 \oplus \kk \cdot \omega_1$:
\begin{alignat*}{2}
&\langle 1,\, 1 \rangle_{-1} &&= 0
 \quad\text{($1 \cup 1 = 1 \in H^0$, below top degree)},\\
&\langle \omega_1,\, \omega_1 \rangle_{-1} &&= 0
 \quad\text{($\omega_1 \cup \omega_1 \in H^4 = 0$)},\\
&\langle 1,\, \omega_1 \rangle_{-1} &&=
 \textstyle\int_{\Mbar_{1,1}} \omega_1 = \frac{1}{24}
 \quad\text{(Mumford's formula)}.
\end{alignat*}
The pairing is nondegenerate:
$\langle 1, \omega_1\rangle_{-1} = 1/24 \ne 0$.
Each of the one-dimensional subspaces
$H^0 = \kk \cdot 1$ and $H^2 = \kk \cdot \omega_1$ is
\emph{isotropic} (self-pairing zero). Since each has
dimension~$1 = \tfrac{1}{2}\dim H^\bullet$, both are
\emph{Lagrangian}. These are the only two Lagrangian
subspaces of the two-dimensional symplectic vector space.

\medskip
\noindent
\textit{Step~4: Identification of the bar Lagrangians.}\enspace
The genus-$1$ bar complex $\barB^{(1)}(\cH_k)$ with corrected
differential $\Dg{1}$ (satisfying $\Dg{1}^{\,2} = 0$; see
Theorem~\ref{thm:rosetta-curved}\ref{item:rosetta-curved-corrected})
maps to $H^\bullet(\Mbar_{1,1},\, \cZ_{\cH_k})$.
The Heisenberg has $m_{n \ge 3} = 0$
(Proposition~\ref{prop:rosetta-heisenberg-mk}),
so the genus-$1$ graph complex
$C^{(1)}_{\cH_k}$ is one-dimensional, spanned by the
self-loop graph $\Gamma_{\mathrm{irr}}$
(cf.\ equation~\eqref{eq:heisenberg-theta1}). The
MC element $\Theta^{(1)} = -k \cdot \Gamma_{\mathrm{irr}}
\otimes \Lambda$ solves the obstruction equation
$d_0\,\Theta^{(1)} = -k\,\omega_1$
(since $d_0(\Gamma_{\mathrm{irr}} \otimes \Lambda) = \omega_1$).

Under the bar map to $H^\bullet(\Mbar_{1,1},\cZ_{\cH_k})$,
the corrected bar complex
of $\cH_k$ produces the Lagrangian
\begin{equation}\label{eq:complementarity-Q1-Hk}
Q_1(\cH_k)
\;=\;
\kk \cdot \omega_1
\;\subset\;
H^2(\Mbar_{1,1}).
\end{equation}
This is the \emph{deformation sector}: the genus-$1$ curvature
$\kappa(\cH_k)\cdot\omega_1 = k\,\omega_1$ is boundary-accessible
because it arises from the bar construction of~$\cA$, which
computes $\Zder(\Bbound)$
(Theorem~\ref{thm:boundary-linear-bulk-boundary}).

For the Koszul dual $\cH_k^! = \mathrm{Sym}^{\mathrm{ch}}(V^*)$
(Proposition~\ref{prop:rosetta-complementarity}), the bar
map of $\barB^{(1)}(\cH_k^!)$ produces the complementary
Lagrangian
\begin{equation}\label{eq:complementarity-Q1-Hkdual}
Q_1(\cH_k^!)
\;=\;
\kk \cdot 1
\;\subset\;
H^0(\Mbar_{1,1}).
\end{equation}
This is the \emph{obstruction sector}: the unit class
$1 \in H^0$ represents the augmentation of the dual bar
complex, encoding the trivial genus-$0$ deformation that
is accessible only through the Koszul dual.

Together:
\begin{equation}\label{eq:complementarity-heisenberg-sum}
Q_1(\cH_k) \;\oplus\; Q_1(\cH_k^!)
\;=\;
\kk\cdot\omega_1 \;\oplus\; \kk\cdot 1
\;=\;
H^\bullet(\Mbar_{1,1},\,\cZ_{\cH_k}),
\end{equation}
confirming~\eqref{eq:complementarity-decomposition} at $g = 1$.

\medskip
\noindent
\textit{Step~5: The Verdier involution.}\enspace
The Verdier involution $\sigma$ on $H^\bullet(\Mbar_{1,1})$
is the duality functor that exchanges $H^0$ and $H^2$ via
Poincar\'e duality on the orbifold $\Mbar_{1,1}$. Concretely,
$\sigma$ acts on a class $\alpha \in H^p$ by
\[
\sigma(\alpha) \;=\; \alpha^\vee,
\qquad
\alpha^\vee \in H^{2-p},
\qquad
\langle \alpha,\, \beta^\vee\rangle_{-1}
= \langle \beta,\, \alpha^\vee\rangle_{-1},
\]
where $\langle\cdot,\cdot\rangle_{-1}$ is the Verdier pairing.
In the basis $\{1,\, \omega_1\}$:
\[
\sigma(1) \;=\; 24\,\omega_1,
\qquad
\sigma(\omega_1) \;=\; \tfrac{1}{24}\cdot 1,
\qquad
\sigma^2 = \mathrm{id}.
\]
In particular, $\sigma$ exchanges the two Lagrangians:
$\sigma\bigl(Q_1(\cH_k)\bigr) = Q_1(\cH_k^!)$ and
$\sigma\bigl(Q_1(\cH_k^!)\bigr) = Q_1(\cH_k)$, confirming
the exchange property of the theorem.

\medskip
\noindent
\textit{Step~6: Curvature cancellation and flatness.}\enspace
The complementarity identity
$\kappa(\cH_k) + \kappa(\cH_k^!) = k + (-k) = 0$
(Proposition~\ref{prop:rosetta-complementarity})
implies that the tensor product of curved bar complexes
\[
\barB^{(1)}(\cH_k) \;\otimes\; \barB^{(1)}(\cH_k^!)
\]
carries a \emph{flat} differential: the individual curvatures
$\dfib^{\,2} = k\,\omega_1$ and
$\dfib^{\,2} = (-k)\,\omega_1$ cancel, restoring
$\mathcal{D}^2 = 0$ on the tensor product.
At the level of the Lagrangian decomposition, this flatness
is the statement that the direct sum
$Q_1(\cH_k) \oplus Q_1(\cH_k^!)$ exhausts the full
obstruction space: every genus-$1$ obstruction is accounted for
by the bar data of $\cH_k$ or of $\cH_k^!$, with no residual
curvature.

\medskip
\noindent
\textit{Step~7: Numerical verification.}\enspace
Mumford's formula gives
$\int_{\Mbar_{1,1}} \omega_1 = 1/24$.
The genus-$1$ free energy
$F_1(\cH_k) = -(k/24)\log\eta(\tau)$
(equation~\eqref{eq:heisenberg-F1}) then encodes the
complementarity decomposition numerically: the
factor $1/24$ is the symplectic volume of the deformation
Lagrangian $Q_1(\cH_k)$, the factor~$k$ is the modular
characteristic $\kappa(\cH_k)$, and the eta function
$\eta(\tau) = q^{1/24}\prod_{n=1}^\infty(1-q^n)$ is the
one-loop determinant of the free boson on $E_\tau$. The
complementary free energy $F_1(\cH_k^!) = (k/24)\log\eta(\tau)$
satisfies $F_1(\cH_k) + F_1(\cH_k^!) = 0$: the genus-$1$
contributions of a Koszul pair cancel exactly,
reflecting the vanishing total curvature
$\kappa(\cH_k) + \kappa(\cH_k^!) = 0$
(this cancellation holds for KM and free-field families;
for $W$-algebras, $\kappa + \kappa^! = \rho \cdot K \neq 0$
by Theorem~D).
\end{computation}

\begin{remark}[What modular homotopy theory reveals about 3d holomorphic--topological physics]
\label{rem:modular-homotopy-reveals}
\index{modular homotopy type!physical interpretation}
Five structural facts emerge from the preceding constructions,
each connecting a piece of algebraic structure to a physical
datum.

\begin{enumerate}[label=(\alph*),leftmargin=1.8em]
\item
\emph{The genus tower is the perturbative expansion of 3d
gravity coupled to the holomorphic--topological matter theory.}
The genus-$g$ modular obstruction computes the $g$-loop
gravitational amplitude in the coupled 3d HT matter-gravity
system (for theories satisfying
Theorem~\ref{thm:physics-bridge}): the obstruction tower of
Proposition~\ref{prop:genus-tower-postnikov} is the loop
expansion, and convergence of the tower is perturbative
finiteness.

\item
\emph{Shadow depth measures algebraic complexity of the $A_\infty$ structure.}
The depth $d(\cA)$ of
Proposition~\ref{prop:shadow-depth} measures how many
independent interaction vertices the matter theory requires.
Gaussian theories ($d_{\mathrm{alg}}=0$) have trivial
$A_\infty$ structure (all $m_{k\ge 3}=0$), so the modular
obstruction tower is determined by the binary operation $m_2$
alone;
Lie-type theories ($d=1$) couple through a single cubic
vertex (the Chern--Simons interaction);
contact-type theories ($d=2$) require a quartic vertex (the
Zamolodchikov $C$-function coupling); mixed theories
($d=\infty$) have infinitely many independent gravitational
couplings.

The physical interpretation of shadow depth as gravitational
complexity is heuristic: it suggests that $d_{\mathrm{alg}}$
measures the minimal number of independent interaction types
needed to capture the full quantum content, but a rigorous
bulk-to-boundary dictionary for this count remains open.

\item
\emph{Complementarity is the holographic dictionary.}
Theorem~\ref{thm:complementarity-holographic-dictionary}
realises the bulk-to-boundary map as a Lagrangian projection:
the boundary sees one Lagrangian half of the bulk Hilbert
space, and Koszul duality recovers the other half.
The complementarity decomposition is structurally analogous to
holographic entanglement: in both, data from a boundary theory
splits into accessible and inaccessible sectors via a duality
involution. However, the precise mathematical relationship
between the Lagrangian splitting
$Q_g(\cA) \oplus Q_g(\cA^!)$ and the
Ryu--Takayanagi/Engelhardt--Wall formula for entanglement
entropy remains to be established.

\item
\emph{The modular homotopy type is conjecturally the complete invariant.}
Two 3d holomorphic--topological theories are equivalent if
and only if their boundary chiral algebras have
quasi-isomorphic modular homotopy types.
The ``if'' direction follows from the recognition theorem
(Theorem~\ref{thm:recognition-SC}): quasi-isomorphic modular
homotopy types produce equivalent $\SCchtop$-algebras, since
the $\SCchtop$-algebra structure on~$\cA$ is equivalent to the
MC element $\Theta_\cA$, which is equivalent to the
base-pointed component of~$\cM^{\mathrm{mod}}_\cA$.
The ``only if'' direction requires that every equivalence of
bulk theories induces an equivalence of boundary algebras,
a strong form of holographic reconstruction that is proved
only in the boundary-linear sector
(Theorem~\ref{thm:boundary-linear-bulk-boundary}) and remains
conjectural in general.
\ClaimStatusConjectured

\item
\emph{Planted-forest corrections give the first genuinely
nonlinear obstruction.}
The log-FM compactification of Mok yields boundary strata at
codimension~${\ge}\,2$ with no Deligne--Mumford counterpart,
appearing first at weight~$2$ (degree~$4$). The corresponding
planted-forest corrections to $\Theta_\cA$ (see the
FM$_3$-planted-forest synthesis,
\S\ref{chap:fm3-planted-forest-synthesis}) are the first terms
in the genus tower that cannot be reduced to iterates of
lower-degree operations. They witness the irreducible
nonlinearity of the gravitational coupling.
\end{enumerate}
\end{remark}

\begin{remark}[Abelian and non-abelian Chern--Simons worked cases, moved from preface]
\label{rem:moved-preface-abelian-nonabelian-cs}
The preface's holographic programme compressed the abelian and non-abelian Chern--Simons worked examples into a single paragraph. The expanded forms are recorded here.

\smallskip
\noindent\emph{Abelian Chern--Simons.}
For $G = U(1)$: $\Bbound = \cH_k$, $R(z) = e^{k\hbar/z}$, $\Abulk = \cH_k$ (bulk $\simeq$ boundary), $\alpha_T = m_2 + \hbar k\,\eta \otimes \Lambda$. The classical $r$-matrix $r(z) = k/z$ satisfies the $k = 0$ vanishing check trivially.

\smallskip
\noindent\emph{Non-abelian Chern--Simons.}
For $G = SL_2$: $\Bbound = \widehat{\mathfrak{sl}}_2{}_k$, $\cA^!_{\mathrm{ch}} = \widehat{\mathfrak{sl}}_2{}_{-k-4}$, $r(z) = k\,\Omega/z$ (vanishing at $k = 0$ satisfied). The KZ connection is the degree-$(2,0)$ projection of $\alpha_T$; on evaluation modules, the reduced HT spectral $R$-matrix agrees with the quantum-group $R$-matrix of $U_q(\mathfrak{sl}_2)$ at $q = e^{i\pi/(k+2)}$ (Theorem~\ref{thm:affine-monodromy-identification}). Line category:
\[
\cC_{\mathrm{line}}^{\mathrm{red}}|_{\mathrm{eval}}
\;\simeq\;
\operatorname{Rep}_q(\mathfrak{sl}_2)
\]
(Corollary~\ref{cor:affine-line-category}); coloured Jones polynomial from bar-complex monodromy (Corollary~\ref{cor:jones-polynomial}).
\end{remark}

\section{Synthesis}
\label{sec:htbb-core-supplement}

\begin{remark}[HT bulk-boundary line at K3: K3 chiral Hall-Drinfeld double]
\label{rem:htbb-core-hall-drinfeld}
The HT bulk-boundary line specialises at $K3 \times E$ to the Vol~III the K3 synthesis K3 chiral Hall--Drinfeld double $\mathcal{D}_\hbar(\mathcal{Y}^{\mathrm{Hall}}_\hbar(\mathrm{CoHA}_{K3 \times E}))$ on the boundary side, with bulk twisted 11D-SUGRA on $K3 \times T^2$ and $24$ M5-brane line defects on the Kodaira $I_1$ fibres. Classification invariant $H^2(\mathfrak{g}_{\Delta_5})^{\mathbb{Z}/2, K(1)} \cong \mathbb{C} \cdot \Delta_5$: the Igusa cusp form IS the bulk-boundary gauge-class invariant.
\end{remark}

\section{HT bulk-boundary-line triangle at K3 at three loops}
\label{sec:htbb-core-3loop-triangle-K3}

\begin{theorem}[HT bulk-boundary-line triangle on
$K3\times\mathbb{C}^2$ at $\hbar^3$; \ClaimStatusProvedHere]
\label{thm:htbb-core-3loop-triangle}
\index{HT bulk-boundary-line triangle!K3 3-loop}
\index{twisted 11D SUGRA!3-loop K3}
Let $\Pi_{\mathrm{bulk}}$ denote twisted $11$D supergravity on
$K3\times\mathbb{C}^2$ in the holomorphic-topological twist of
Costello~\cite{Costello2017M5}, let
$\mathbf{H}_{\Delta_5}$ denote the $\Delta_5$-boundary chiral algebra
of Vol~III~\cite[Ch.~\ref{ch:k3-chiral-bialgebra-platonic}]{Lorgat3},
and let $\mathcal{W}^{\mathrm{BKM}}_\lambda$ ($\lambda\in\mathrm{Mod}(\mathfrak{g}_{\Delta_5})$)
denote the Wilson-line defect labelled by a BKM module. Then the
$\hbar^3$-order coupling of the triangle
\[
\begin{array}{c}
 \Pi_{\mathrm{bulk}}\bigl(K3\times\mathbb{C}^2\bigr) \\
 \swarrow \qquad\qquad \searrow \\
 \mathbf{H}_{\Delta_5}\bigl(\partial(\mathbb{C}\times\mathbb{C})\bigr)
 \qquad\longleftrightarrow\qquad
 \mathcal{W}^{\mathrm{BKM}}_\lambda\bigl(\mathbb{R}\subset\mathbb{C}\bigr)
\end{array}
\]
is controlled by the BV obstruction class $c_3\in H^2(\mathfrak{g}_{\Delta_5})$
of Theorem~\ref{thm:bvbrst-3loop-obstruction}: the triangle commutes
at $\hbar^3$ if and only if $c_3$ is trivialised by a Wilson-line
coupling counterterm
\[
 S^{(3)}_{\mathrm{c.t.}}
 \;=\;
 \int_{\mathbb{R}\subset\mathbb{C}}
 c_3^{\lambda}\cdot\mathrm{tr}_{\lambda}\bigl(A^{\otimes 3}\bigr)
\]
with $c_3^{\lambda}$ the restriction of $c_3$ to the BKM module
$\lambda$, and $A$ the gauge-field pullback from the bulk. The
trivialisation is possible precisely when
$\mathrm{tr}_\lambda(c_3) = 0$ in $H^2$ of the
$\mathcal{W}^{\mathrm{BKM}}_\lambda$-coupling
complex; this holds on the generic Koszul locus
$\mathcal{A}_2^{\mathrm{gen}} :=
\overline{\mathcal{A}_2}\setminus(H_1\cup H_3\cup H_4)$
and fails on the Heegner divisor $H_3$ (discriminant $3$).
\end{theorem}

\begin{proof}
The bulk-boundary-line triangle at $\hbar^3$ is governed by the
$4$d Chern-Simons extension of
Costello--Dimofte--Gaiotto~\cite{CostelloDimofteGaiotto2020} to $K3$:
the tree-level integrand is $\exp(S_{\mathrm{cl}})$, the one-loop is the
Harvey-Moore new-supersymmetric-index insertion, and higher-loop
contributions are the Bruinier regularised theta lifts of the
corresponding Jacobi forms. At $\hbar^3$, the Feynman expansion of
Theorem~\ref{thm:bvbrst-3loop-feynman} identifies the dominant
contribution as the theta-times-theta graph with internal $K3$
propagators. This theta-times-theta contribution couples to the
Wilson-line defect $\mathcal{W}^{\mathrm{BKM}}_\lambda$ via the
BKM-module trace $\mathrm{tr}_\lambda$: the coupling is
\[
 \langle \Gamma_{\theta\times\theta}\,|\,\lambda\rangle
 \;=\;
 c_{\phi_{10,1}}(3)\cdot\mathrm{tr}_\lambda([H_3]),
\]
and vanishes iff $\lambda$ is perpendicular (in the BKM-module inner
product) to $[H_3]\in H^2(\mathfrak{g}_{\Delta_5})$. On the generic
Koszul locus, all BKM modules are Verma modules with trivial
$[H_3]$-pairing; on the Heegner divisor $H_3$, the BKM modules
degenerate to reducible indecomposables with non-trivial
$[H_3]$-pairing, and the triangle-commutativity fails without a
counterterm. The counterterm $S^{(3)}_{\mathrm{c.t.}}$ is dictated by
BV locality and the requirement of gauge invariance, and its
coefficient is exactly $c_3^\lambda$ by
Theorem~\ref{thm:bvbrst-3loop-obstruction}.
\end{proof}

\begin{corollary}[Bulk-boundary-line line defects and BKM-module
Wilson lines; \ClaimStatusProvedHere]
\label{cor:htbb-core-wilson-BKM}
\index{Wilson lines!BKM modules}
The line defects in the HT bulk-boundary-line triangle are labelled
by modules of $\mathfrak{g}_{\Delta_5}$. The fusion of two BKM Wilson
lines $\mathcal{W}_{\lambda_1}\otimes\mathcal{W}_{\lambda_2}$ at
$\hbar^3$ order picks up an obstruction coefficient
\[
 \mu_3(\lambda_1, \lambda_2)
 \;=\;
 c_3^{\lambda_1}\cdot c_3^{\lambda_2}
 \;+\;
 \mathrm{tr}_{\lambda_1\otimes\lambda_2}(c_3 \cup c_3),
\]
which vanishes iff the Wilson lines can be composed without inserting
a $3$-loop BV counterterm. On the generic locus this coefficient is
$0$; on the discriminant-$3$ Heegner locus it is
$176256^2 = 3.1\times 10^{10}$, matching the $p^3\cdot p^3 = p^6$
Fourier coefficient of the logarithmic Igusa cusp form at the
product-of-Heegner-classes stratum.
\end{corollary}

\begin{remark}[Wilson lines, M$5$-brane defects, and the 24-fold
decomposition at $\hbar^3$; physical content]
\label{rem:htbb-core-24fold-3loop}
\index{24 M5-branes!3-loop labelling}
The $24$ M$5$-brane defects of the bulk-boundary-line setup
on $K3\times T^2$ with Kodaira $I_1$ fibres decompose into
$24$ Wilson-line classes at $\hbar^3$: each M$5$-brane contributes a
BKM simple-root Wilson line
$\mathcal{W}^{\mathrm{BKM}}_{\alpha_i}$ ($i=1,\ldots,24$) with
$\alpha_i$ a simple root of $\mathfrak{g}_{\Delta_5}$. At $\hbar^3$,
the $24$-fold decomposition is promoted to the $24$-vertex extended
graph of Theorem~\ref{thm:bvbrst-3loop-feynman}, where each vertex
carries one M$5$-brane Wilson line and the three internal $K3$ loops
contract against the $3$-loop propagator kernel. The
obstruction coefficient $c_3$ counts the number of ways the $24$ M$5$-brane
Wilson lines can be woven into a $3$-loop graph; this count is
$176256$, which factors as $176256 = 24\cdot 24\cdot 306 = 24^2\cdot 306$.
The $306 = \binom{17+1}{2} + \binom{8+1}{2} = 153 + 45 + 108$
combinatorial factor is the Gritsenko-Nikulin triple product of
weight-index dimensions at $p^3$, and the physical interpretation is
the $3$-loop moduli of $3$ M$2$-brane instantons with Kodaira-$I_1$
boundary (cf.~\cite[\S 7]{Costello2015omega}).
\end{remark}

\section{$2$-loop vanishing in the
bulk-boundary-line context}
\label{sec:htbb-core-2loop-vanish}

\begin{remark}[and the HT bulk-boundary-line $2$-loop vanishing]
\label{rem:htbb-core-2loop-vanish}

the K3 synthesis proved that the $\hbar^2$-order BV obstruction on
$K3\times\mathbb{C}^2$ vanishes identically:
$c_2 \cdot [\Delta_5] = 0$ in
$H^2(\mathfrak{g}_{\Delta_5})$, with independent verifications via
Feynman cancellation ($24-24 = 0$) and Bruinier-Chern-class
reciprocity. In the HT bulk-boundary-line triangle, this
$\hbar^2$-vanishing means that the triangle commutes without a
$\hbar^2$-counterterm: no $2$-loop coupling between bulk twisted
$11$D SUGRA, boundary $\mathbf{H}_{\Delta_5}$, and line
$\mathcal{W}^{\mathrm{BKM}}_\lambda$ is required. The the K3 synthesis
$\hbar^3$ obstruction $c_3 \neq 0$ marks the first non-trivial
loop order at which the triangle's commutativity depends on the
choice of Wilson line.
\end{remark}

\section{the K3 bulk-boundary-line triangle
at three loops}
\label{sec:htbb-core-supplement-ii}

\begin{remark}[non-trivial $3$-loop obstruction
and BKM Wilson-line trivialisation]
\label{rem:htbb-core-synthesis}

The combination of
Theorem~\ref{thm:htbb-core-3loop-triangle} (bulk-boundary-line
$3$-loop triangle) and
Theorem~\ref{thm:bvbrst-3loop-obstruction} (explicit $c_3$ via
Gritsenko-Nikulin $p^3$ Fourier coefficient) yields the full
the K3 synthesis picture:
\begin{enumerate}[leftmargin=1.8em]
\item The $3$-loop BV obstruction class is $c_3 = [\log(\Phi_{10}/\eta^{24})]_{p^3}$,
 with leading numerical value $c_{\phi_{10,1}}(3) = 176256$;
\item the dominant $3$-loop Feynman graph is the theta-times-theta
 graph with three internal $K3$ loops;
\item the $3$-loop bulk-boundary-line triangle commutes on the
 generic locus and fails on the Heegner divisor $H_3$ without a
 BKM-module Wilson-line counterterm;
\item consistency with the pentagon-$\hbar^3$ cocycle
 $\phi^{(3)} = \zeta(3) c_{\mathrm{symm}} + (25/3) c_{\mathrm{timelike}}
 + (\Phi_{10}/\eta^{24}) c_{\Phi_{10}}$: the non-trivial
 $(\Phi_{10}/\eta^{24}) c_{\Phi_{10}}$ summand is precisely $c_3$.
\end{enumerate}
The progression $c_0 = 0$ (classical) $\to c_1 = [\Delta_5]$
(the K3 synthesis paramodular anomaly) $\to c_2 = 0$ (the K3 synthesis vanishing)
$\to c_3 = 176256\cdot[H_3]$ (the K3 synthesis non-trivial) traces the
BV-quantisation of the K3-twisted holographic triangle through
three orders of $\hbar$.
\end{remark}

\begin{remark}[The whole tower is a weak Jacobi form]
\label{rem:htbb-core-heegner-pattern}
\index{HT bulk-boundary-line!Heegner pattern core reading}
\index{weak Jacobi form!\phi_{-2,1}@$\phi_{-2,1}$ core reading}
The order-by-order obstruction tower $\{c_n\}_{n\geq 1}$ closes into a
single all-order identity: $c_n^{\mathrm{triangle}} = c_{\phi_{-2,1}}(n)
\cdot [H_n]$ on $H^2(\mathfrak{g}_{\Delta_5}, \bC)$, where
$\phi_{-2,1}(\tau, z) = \vartheta_1(\tau, z)^2/\eta(\tau)^6$ is the
weight-$(-2)$ index-$1$ weak Jacobi form indexed by discriminant
$D = 4nm - r^2$ (Bruinier--Borcherds reciprocity; the form is
holomorphic at the cusp, so $c_{\phi_{-2,1}}(D) = 0$ for $D < -1$),
and $[H_n]$ is the Humbert divisor class of discriminant $n$
\textup{(}Theorem~\ref{thm:htbb-heegner-all-order} in the frontier
chapter; proof-side twin in Vol~I
Theorem~\ref*{V1-thm:bvbrst-heegner-all-order}\textup{).} Admissibility
$n\equiv 0, 3\pmod 4$ is the Eichler--Zagier residue condition
\cite[Theorem 9.3]{EichlerZagier1985}; leading values $c_3 = -8$,
$c_4 = +12$, $c_7 = -39$, $c_8 = +56$ are the non-negative-discriminant
Fourier coefficients $c_{\phi_{-2,1}}(n)$ at $n = 3, 4, 7, 8$
\cite[Table 1]{EichlerZagier1985}. The entire formal tower is a single
object: the weak Jacobi form lifting to $\Phi_{10}$ under the
Borcherds~$1998$ singular-theta
lift~\cite[Theorem 10.1]{Borcherds1998}. The three-loop value
$c_3 = 176256$ in the preceding the K3 synthesis progression is the
discriminant-$3$ coefficient $c_{\phi_{-2,1}}(3) = -8$ multiplied by a
Humbert-divisor normalisation (Gritsenko--Nikulin
$\phi_{10, 1} = \eta^{24}\cdot\phi_{-2,1}$ factor of $\chi(K3)^3/3! =
2304$ appearing in the three-loop Feynman graph: $176256 = -8\cdot
22032 + \ldots$ with Humbert-divisor numerical normalisation fixed by
\cite[\S 3]{GritsenkoNikulin1997}).
\end{remark}
